
\chapter{群导论}

\section{基本公理与例子}

本节引入第一部分将要研究的基本代数结构，并给出一些例子。

\noindent\textbf{定义.}
\begin{enumerate}[label=(\arabic*)]
  \item 集合 $G$ 上的一个\emph{二元运算} $*$ 是一个函数 $*:G\times G\to G$。对任意 $a,b\in G$，我们把 $*(a,b)$ 写作 $a*b$。
  \item 若集合 $G$ 上的二元运算 $*$ 满足对所有 $a,b,c\in G$，都有
  \[
    a*(b*c)=(a*b)*c,
  \]
  则称 $*$ 是\emph{结合的}。
  \item 若 $*$ 是集合 $G$ 上的二元运算，若 $G$ 中元素 $a,b$ 满足 $a*b=b*a$，则称 $a$ 与 $b$ \emph{交换}。若对所有 $a,b\in G$ 都有 $a*b=b*a$，则称 $*$（或称 $G$）是\emph{交换的}。
\end{enumerate}

\paragraph{例}
\begin{enumerate}[label=(\arabic*)]
  \item $+$（通常加法）是 $\ZZ$ 上的交换二元运算（在 $\QQ$、$\RR$ 或 $\CC$ 上亦然）。
  \item $\times$（通常乘法）是 $\ZZ$ 上的交换二元运算（在 $\QQ$、$\RR$ 或 $\CC$ 上亦然）。
  \item $-$（通常减法）是 $\ZZ$ 上的非交换二元运算，其中 $-(a,b)=a-b$。映射 $a\mapsto -a$ 不是二元运算（不是二元的）。
  \item $-$ 不是 $\ZZ^+$ 上的二元运算（也不是 $\QQ^+$ 或 $\RR^+$ 上的二元运算），因为若 $a,b\in\ZZ^+$ 且 $a<b$，则 $a-b\notin\ZZ^+$；也就是说，$-$ 并不把 $\ZZ^+\times\ZZ^+$ 映入 $\ZZ^+$。
  \item 三维空间 $\RR^3$ 中两个向量的叉积是一个既不结合也不交换的二元运算。
\end{enumerate}

设 $*$ 是集合 $G$ 上的二元运算，且 $H$ 是 $G$ 的子集。若 $*$ 到 $H$ 的限制是 $H$ 上的二元运算，即对所有 $a,b\in H$ 都有 $a*b\in H$，则称 $H$ 对 $*$ \emph{封闭}。注意，如果 $*$ 是 $G$ 上的结合（相应地，交换）二元运算，并且 $*$ 限制到 $G$ 的某个子集 $H$ 上是 $H$ 的二元运算，那么 $*$ 在 $H$ 上也自动是结合的（相应地，交换的）。

\noindent\textbf{定义.}
\begin{enumerate}[label=(\arabic*)]
  \item \emph{群}是一个有序对 $(G,*)$，其中 $G$ 是一个集合，$*$ 是 $G$ 上的二元运算，并满足以下公理：
  \begin{enumerate}[label=(\roman*)]
    \item 对所有 $a,b,c\in G$，$(a*b)*c=a*(b*c)$，即 $*$ 是结合的；
    \item 存在 $G$ 中的元素 $e$，称为 $G$ 的\emph{单位元}，使得对所有 $a\in G$，都有 $a*e=e*a=a$；
    \item 对每个 $a\in G$，存在 $G$ 中的元素 $a^{-1}$，称为 $a$ 的\emph{逆元}，使得 $a*a^{-1}=a^{-1}*a=e$。
  \end{enumerate}
  \item 若群 $(G,*)$ 满足对所有 $a,b\in G$ 都有 $a*b=b*a$，则称它是\emph{阿贝尔的}（或\emph{交换的}）。
\end{enumerate}

我们很快会不那么形式化：若 $(G,*)$ 是一个群，就说 $G$ 在 $*$ 下是一个群（当运算 $*$ 从上下文中清楚时，干脆说 $G$ 是一个群）。此外，若 $G$ 还是一个有限集，则称 $G$ 是\emph{有限群}。注意，公理 (ii) 保证群总是非空的。

\paragraph{例}
\begin{enumerate}[label=(\arabic*)]
  \item $\ZZ$、$\QQ$、$\RR$ 和 $\CC$ 在 $+$ 下是群，其中 $e=0$，并且对所有 $a$，$a^{-1}=-a$。
  \item $\QQ-\{0\}$、$\RR-\{0\}$、$\CC-\{0\}$、$\QQ^+$、$\RR^+$ 在 $\times$ 下是群，其中 $e=1$，并且对所有 $a$，$a^{-1}=1/a$。但是要注意，$\ZZ-\{0\}$ 在 $\times$ 下不是群，因为虽然 $\times$ 是 $\ZZ-\{0\}$ 上的结合二元运算，但元素 $2$（例如）在 $\ZZ-\{0\}$ 中没有逆元。
\end{enumerate}

我们在这些熟悉的例子中略过了结合律成立这一事实。对 $\ZZ$ 的加法而言，它是自然数加法结合公理的结果。$\QQ$ 上加法的结合律由 $\ZZ$ 的结合律推出——稍后在从 $\ZZ$ 严格构造 $\QQ$ 时会给出证明概要（参见第 7.5 节）。$\RR$ 以及进而 $\CC$ 上加法的结合律，则在初等分析课程中通过把 $\RR$ 构造为 $\QQ$ 的完备化来证明；最终，结合律仍是 $\ZZ$ 结合律的结果。乘法的结合公理也可通过类似的发展建立，先从 $\ZZ$ 开始。由于 $\RR$ 和 $\CC$ 在本书中主要用于说明性目的，而且我们不会从 $\QQ$ 构造 $\RR$（虽然会从 $\RR$ 构造 $\CC$），所以我们把 $\RR$ 与 $\CC$ 在 $+$ 和 $\times$ 下的结合律作为已知事实。

\paragraph{例（续）}
\begin{enumerate}[label=(\arabic*),start=3]
  \item 向量空间 $V$ 的公理包括规定 $(V,+)$ 是阿贝尔群的那些公理（运算 $+$ 称为向量加法）。因此，任何向量空间，例如 $\RR^n$，特别地都是一个加法群。
  \item 对 $n\in\ZZ^+$，$\ZZ/n\ZZ$ 在第 0 章所述的剩余类加法 $+$ 下是阿贝尔群。我们将在第 3 章（在更一般的语境中）证明这个二元运算 $+$ 是良定且结合的；目前先承认这一点。这个群的单位元是元素 $\overline{0}$，对每个 $\overline{a}\in\ZZ/n\ZZ$，其逆元是 $\overline{-a}$。今后当我们谈到群 $\ZZ/n\ZZ$ 时，默认其群运算是模 $n$ 的类加法。
  \item 对 $n\in\ZZ^+$，集合 $(\ZZ/n\ZZ)^\times$ 由所有在模 $n$ 下有乘法逆元的等价类 $\overline{a}$ 组成；它在第 0 章所述的剩余类乘法下是阿贝尔群。同样，我们暂且承认该运算是良定且结合的。这个群的单位元是元素 $\overline{1}$，并且按 $(\ZZ/n\ZZ)^\times$ 的定义，每个元素都有乘法逆元。今后当我们谈到群 $(\ZZ/n\ZZ)^\times$ 时，默认其群运算是模 $n$ 的类乘法。
  \item 若 $(A,*)$ 和 $(B,\circ)$ 是群，则可以构造一个新群 $A\times B$，称为它们的\emph{直积}，其元素是笛卡儿积
  \[
    A\times B=\{(a,b)\mid a\in A,\ b\in B\},
  \]
  中的元素，运算按分量定义：
  \[
    (a_1,b_1)(a_2,b_2)=(a_1*a_2,\ b_1\circ b_2).
  \]
  例如，若取 $A=B=\RR$（二者的运算都是加法），则 $\RR\times\RR$ 就是熟悉的欧氏平面。两个群的直积仍是群的证明留作一个直接的练习（以后会做）：$A\times B$ 中每条群公理的成立，都是 $A$ 与 $B$ 中相应公理成立以及 $A\times B$ 中运算按分量定义的结果。
\end{enumerate}

不要混淆群 $\ZZ/n\ZZ$（加法下）与 $(\ZZ/n\ZZ)^\times$（乘法下），即使后者是前者的子集；上标 $\times$ 将总表示运算是乘法。

在继续更复杂的例子之前，我们先证明两个基本结果，它们尤其使我们可以谈论单位元和元素的逆元。

\medskip
\noindent\textbf{命题 1.} 若 $G$ 在运算 $*$ 下是群，则
\begin{enumerate}[label=(\arabic*)]
  \item $G$ 的单位元唯一；
  \item 对每个 $a\in G$，$a^{-1}$ 唯一确定；
  \item 对所有 $a\in G$，$(a^{-1})^{-1}=a$；
  \item $(a*b)^{-1}=(b^{-1})*(a^{-1})$；
  \item 对任意 $a_1,a_2,\ldots,a_n\in G$，$a_1*a_2*\cdots*a_n$ 的值与该表达式如何加括号无关（这称为\emph{广义结合律}）。
\end{enumerate}

\begin{proof}
(1) 若 $f$ 和 $g$ 都是单位元，则由群定义的公理 (ii)，$f*g=f$（取 $a=f$、$e=g$）。同一公理也给出 $f*g=g$（取 $a=g$、$e=f$）。故 $f=g$，单位元唯一。

(2) 假设 $b$ 和 $c$ 都是 $a$ 的逆元，令 $e$ 为 $G$ 的单位元。由公理 (iii)，$a*b=e$ 且 $c*a=e$。于是
\begin{align*}
  c &= c*e &&\text{（$e$ 的定义——公理 (ii)）}\\
    &= c*(a*b) &&\text{（因为 $e=a*b$）}\\
    &= (c*a)*b &&\text{（结合律）}\\
    &= e*b &&\text{（因为 $e=c*a$）}\\
    &= b &&\text{（公理 (ii)）。}
\end{align*}

(3) 证明 $(a^{-1})^{-1}=a$，正是要证明 $a$ 是 $a^{-1}$ 的逆元（由于第 (2) 部分，$a^{-1}$ 的逆元唯一）。把 $a$ 与 $a^{-1}$ 的角色互换来看 $a^{-1}$ 的定义，便知 $a$ 满足 $a^{-1}$ 的逆元的定义性质，所以 $a$ 是 $a^{-1}$ 的逆元。

(4) 令 $c=(a*b)^{-1}$，则按 $c$ 的定义，$(a*b)*c=e$。由结合律，
\[
  a*(b*c)=e.
\]
左乘 $a^{-1}$ 得
\[
  a^{-1}*(a*(b*c))=a^{-1}*e.
\]
左边用结合律、右边用 $e$ 的定义，得
\[
  (a^{-1}*a)*(b*c)=a^{-1},
\]
因此
\[
  e*(b*c)=a^{-1},
\]
从而
\[
  b*c=a^{-1}.
\]
现在左乘 $b^{-1}$ 并同样化简：
\begin{align*}
  b^{-1}*(b*c)&=b^{-1}*a^{-1},\
  (b^{-1}*b)*c&=b^{-1}*a^{-1},\
  e*c&=b^{-1}*a^{-1},\
  c&=b^{-1}*a^{-1},
\end{align*}
这就是所要证明的。

(5) 这是一个用 $n$ 作归纳的很好的练习。先证明 $n=1,2,3$ 的情形。然后假设对任意 $k<n$，$k$ 个元素 $b_1*b_2*\cdots*b_k$ 的任意加括号方式都能在不改变乘积值的情况下化为形如
\[
  b_1*(b_2*(b_3*(\cdots*b_k))\cdots)
\]
的表达式。任意给 $a_1*a_2*\cdots*a_n$ 加括号，必然把它分成两个子乘积，比如
\[
  (a_1*a_2*\cdots*a_k)*(a_{k+1}*a_{k+2}*\cdots*a_n),
\]
每个子乘积又以某种方式加括号。对这两个子乘积分别应用归纳假设，最后把结果化为
\[
  a_1*(a_2*(a_3*(\cdots*a_n))\cdots),
\]
即可完成归纳。
\end{proof}

注意，在命题 1 的整个证明中，我们都小心地没有改变任何乘积的顺序（除非公理 (ii) 和 (iii) 允许这样做），因为 $G$ 可能是非阿贝尔群。

\paragraph{记号}
\begin{enumerate}[label=(\arabic*)]
  \item 对抽象群 $G$ 来说，在计算中不断写运算 $*$ 是很麻烦的。今后（除非必要）我们的抽象群 $G,H$ 等都把运算写作 $\cdot$，并把 $a\cdot b$ 写作 $ab$。鉴于广义结合律，三个或更多群元素的乘积不再加括号（虽然运算仍然是二元运算）。最后，对抽象群 $G$（运算为 $\cdot$），我们把 $G$ 的单位元记为 $1$。
  \item 对任意群 $G$（省略运算 $\cdot$）和 $x\in G$，以及 $n\in\ZZ^+$，由于乘积 $xx\cdots x$（$n$ 项）与加括号方式无关，记作 $x^n$。把 $x^{-1}x^{-1}\cdots x^{-1}$（$n$ 项）记作 $x^{-n}$。令 $x^0=1$，即 $G$ 的单位元。
\end{enumerate}

这个新记号非常简洁。当然，处理具体群时，我们仍使用自然的（给定的）运算。例如，当运算是 $+$ 时，单位元记为 $0$；任一元素 $a$ 的逆元 $a^{-1}$ 写作 $-a$；$a+a+\cdots+a$（$n>0$ 项）写作 $na$；$-a-a-\cdots-a$（$n$ 项）写作 $-na$；并且 $0a=0$。

\medskip
\noindent\textbf{命题 2.} 设 $G$ 是群，$a,b\in G$。方程 $ax=b$ 与 $ya=b$ 在 $G$ 中分别有唯一解 $x,y$。特别地，$G$ 中左右消去律成立，即
\begin{enumerate}[label=(\arabic*)]
  \item 若 $au=av$，则 $u=v$；
  \item 若 $ub=vb$，则 $u=v$。
\end{enumerate}

\begin{proof}
通过左乘 $a^{-1}$ 并化简，可解出 $ax=b$ 的解为 $x=a^{-1}b$。$x$ 的唯一性来自 $a^{-1}$ 的唯一性。类似地，若 $ya=b$，则 $y=ba^{-1}$。若 $au=av$，左乘 $a^{-1}$ 并化简得 $u=v$。右消去律同理。
\end{proof}

命题 2 的一个后果是：若 $a$ 是 $G$ 的任一元素，且对某个 $b\in G$ 有 $ab=e$ 或 $ba=e$，则 $b=a^{-1}$；也就是说，我们不必检验两个等式都成立。另外，若对某个 $b\in G$ 有 $ab=a$（或 $ba=a$），则 $b$ 必为 $G$ 的单位元；也就是说，我们不必检验对所有 $x\in G$ 都有 $bx=xb=x$。

\noindent\textbf{定义.} 设 $G$ 是群，$x\in G$。若存在正整数 $n$ 使 $x^n=1$，则把满足此式的最小正整数 $n$ 称为 $x$ 的\emph{阶}，并记为 $|x|$；此时称 $x$ 的阶为 $n$。若 $x$ 的任一正幂都不是单位元，则定义 $x$ 的阶为无穷，并称 $x$ 是\emph{无限阶}的。

元素 $x$ 的阶的符号不应与绝对值符号混淆（当 $G\subseteq\RR$ 时我们会小心区分）。用同一个符号表示一个元素的阶和一个集合的基数（或阶）似乎不太谨慎；不过我们将看到，群中元素的阶等于其所有（不同）幂组成的集合的基数，因此“阶”这个词的两种用法是自然相关的。

\paragraph{例}
\begin{enumerate}[label=(\arabic*)]
  \item 群中的元素阶为 $1$ 当且仅当它是单位元。
  \item 在加法群 $\ZZ$、$\QQ$、$\RR$ 或 $\CC$ 中，每个非零（即非单位）元素都是无限阶的。
  \item 在乘法群 $\RR-\{0\}$ 或 $\QQ-\{0\}$ 中，元素 $-1$ 的阶为 $2$，其他所有非单位元素都是无限阶的。
  \item 在加法群 $\ZZ/9\ZZ$ 中，元素 $\overline{6}$ 的阶为 $3$，因为 $\overline{6}\neq\overline{0}$，$\overline{6}+\overline{6}=\overline{12}=\overline{3}\neq\overline{0}$，但 $\overline{6}+\overline{6}+\overline{6}=\overline{18}=\overline{0}$，即这个群的单位元。回忆，在加法群中，一个元素的幂就是该元素的整数倍。类似地，元素 $\overline{5}$ 的阶为 $9$，因为 $45$ 是能被 $9$ 整除的 $5$ 的最小正倍数。
  \item 在乘法群 $(\ZZ/7\ZZ)^\times$ 中，元素 $\overline{2}$ 的幂为 $\overline{2},\overline{4},\overline{8}=\overline{1}$，其中 $\overline{1}$ 是这个群的单位元，所以 $\overline{2}$ 的阶为 $3$。类似地，元素 $\overline{3}$ 的阶为 $6$，因为 $3^6$ 是同余于 $1\pmod 7$ 的 $3$ 的最小正幂。
\end{enumerate}

\noindent\textbf{定义.} 设 $G=\{g_1,g_2,\ldots,g_n\}$ 是有限群，且 $g_1=1$。$G$ 的\emph{乘法表}或\emph{群表}是一个 $n\times n$ 矩阵，其第 $i,j$ 项为群元素 $g_ig_j$。

对有限群来说，乘法表在某种意义上包含了关于群的全部信息。然而，从计算角度看，它是一个笨重的对象（大小为群阶的平方）；从视觉角度看，它也不是判断群性质的好工具。可以把群表类比为一张记录全国任意两座城市间距离的表。这样的表有用，并且本质上记录了所有距离关系；但一幅地图（更好的是标出所有距离的地图）显然更便于使用。我们对群论（尤其是有限群）最初的展开，部分目的正是寻找一种更概念化的方式来观察群的内部结构。

\subsection*{练习}

设 $G$ 是一个群。
\begin{enumerate}[label=\arabic*.]
  \item 判断下列二元运算哪些是结合的：
  \begin{enumerate}[label=(\alph*)]
    \item $\ZZ$ 上由 $a*b=a-b$ 定义的运算 $*$；
    \item $\RR$ 上由 $a*b=a+b+ab$ 定义的运算 $*$；
    \item $\QQ$ 上由 $a*b=(a+b)/5$ 定义的运算 $*$；
    \item $\ZZ\times\ZZ$ 上由 $(a,b)*(c,d)=(ad+bc,bd)$ 定义的运算 $*$；
    \item $\QQ-\{0\}$ 上由 $a*b=a/b$ 定义的运算 $*$。
  \end{enumerate}
  \item 判断前一题中的哪些二元运算是交换的。
  \item 证明 $\ZZ/n\ZZ$ 中剩余类加法是结合的（可假设它是良定的）。
  \item 证明 $\ZZ/n\ZZ$ 中剩余类乘法是结合的（可假设它是良定的）。
  \item 证明对所有 $n>1$，$\ZZ/n\ZZ$ 在剩余类乘法下不是群。
  \item 判断下列集合在加法下哪些是群：
  \begin{enumerate}[label=(\alph*)]
    \item 最简形式下分母为奇数的有理数集合（包括 $0=0/1$）；
    \item 最简形式下分母为偶数的有理数集合（包括 $0=0/1$）；
    \item 绝对值 $<1$ 的有理数集合；
    \item 绝对值 $\geq 1$ 的有理数集合再加上 $0$；
    \item 分母等于 $1$ 或 $2$ 的有理数集合；
    \item 分母等于 $1$、$2$ 或 $3$ 的有理数集合。
  \end{enumerate}
  \item 令 $G=\{x\in\RR\mid 0\leq x<1\}$，对 $x,y\in G$，令 $x*y$ 为 $x+y$ 的小数部分（即 $x*y=x+y-[x+y]$，其中 $[a]$ 是不超过 $a$ 的最大整数）。证明 $*$ 是 $G$ 上良定的二元运算，并且 $G$ 在 $*$ 下是阿贝尔群（称为模 $1$ 的实数）。
  \item 令 $G=\{z\in\CC\mid z^n=1\text{ 对某个 }n\in\ZZ^+\}$。
  \begin{enumerate}[label=(\alph*)]
    \item 证明 $G$ 在乘法下是群（称为 $\CC$ 中的单位根群）。
    \item 证明 $G$ 在加法下不是群。
  \end{enumerate}
  \item 令 $G=\{a+b\sqrt2\in\RR\mid a,b\in\QQ\}$。
  \begin{enumerate}[label=(\alph*)]
    \item 证明 $G$ 在加法下是群。
    \item 证明 $G$ 的非零元素在乘法下是群。［“有理化分母”以求乘法逆元。］
  \end{enumerate}
  \item 证明有限群是阿贝尔群，当且仅当它的群表是对称矩阵。
  \item 求加法群 $\ZZ/12\ZZ$ 中每个元素的阶。
  \item 求乘法群 $(\ZZ/12\ZZ)^\times$ 中下列元素的阶：$\overline{1},\overline{-1},\overline{5},\overline{7},\overline{-5},\overline{-7}$。
  \item 求加法群 $\ZZ/36\ZZ$ 中下列元素的阶：$\overline{1},\overline{2},\overline{6},\overline{9},\overline{10},\overline{12},\overline{-1},\overline{-10},\overline{-18}$。
  \item 求乘法群 $(\ZZ/36\ZZ)^\times$ 中下列元素的阶：$\overline{1},\overline{-1},\overline{5},\overline{13},\overline{-13},\overline{17}$。
  \item 证明 $(a_1a_2\cdots a_n)^{-1}=a_n^{-1}a_{n-1}^{-1}\cdots a_1^{-1}$，对所有 $a_1,a_2,\ldots,a_n\in G$ 成立。
  \item 设 $x\in G$。证明 $x^2=1$ 当且仅当 $|x|$ 为 $1$ 或 $2$。
  \item 设 $x\in G$。证明若 $|x|=n$，其中 $n$ 为正整数，则 $x^{-1}=x^{n-1}$。
  \item 设 $x,y\in G$。证明 $xy=yx$ 当且仅当 $y^{-1}xy=x$，当且仅当 $x^{-1}y^{-1}xy=1$。
  \item 设 $x\in G$，$a,b\in\ZZ^+$。
  \begin{enumerate}[label=(\alph*)]
    \item 证明 $x^{a+b}=x^ax^b$ 且 $(x^a)^b=x^{ab}$。
    \item 证明 $(x^a)^{-1}=x^{-a}$。
    \item 对任意整数 $a,b$（正、负或零）建立 (a) 的结论。
  \end{enumerate}
  \item 对 $G$ 中元素 $x$，证明 $x$ 与 $x^{-1}$ 有相同的阶。
  \item 设 $G$ 是有限群，$x\in G$ 的阶为 $n$。证明若 $n$ 为奇数，则对某个 $k$ 有 $x=(x^2)^k$。
  \item 若 $x,g\in G$，证明 $|x|=|g^{-1}xg|$。推出对所有 $a,b\in G$，有 $|ab|=|ba|$。
  \item 假设 $x\in G$ 且 $|x|=n<\infty$。若 $n=st$，其中 $s,t$ 为正整数，证明 $|x^s|=t$。
  \item 若 $a$ 与 $b$ 是 $G$ 中可交换的元素，证明对所有 $n\in\ZZ$，$(ab)^n=a^nb^n$。［先对正 $n$ 用归纳法。］
  \item 证明若对所有 $x\in G$ 都有 $x^2=1$，则 $G$ 是阿贝尔群。
  \item 假设 $H$ 是 $(G,*)$ 的非空子集，且对 $G$ 的二元运算封闭并且对取逆封闭，即对所有 $h,k\in H$，有 $hk\in H$ 且 $h^{-1}\in H$。证明 $H$ 在限制到 $H$ 的运算 $*$ 下是群（这样的子集 $H$ 称为 $G$ 的\emph{子群}）。
  \item 证明若 $x$ 是群 $G$ 的一个元素，则 $\{x^n\mid n\in\ZZ\}$ 是 $G$ 的子群（参见前一题），称为由 $x$ 生成的\emph{循环子群}。
  \item 设 $(A,*)$ 与 $(B,\circ)$ 是群，$A\times B$ 是它们的直积（如例 6 定义）。验证 $A\times B$ 的所有群公理：
  \begin{enumerate}[label=(\alph*)]
    \item 证明结合律成立：对所有 $(a_i,b_i)\in A\times B$，$i=1,2,3$，
    \[
      (a_1,b_1)[(a_2,b_2)(a_3,b_3)]=[(a_1,b_1)(a_2,b_2)](a_3,b_3);
    \]
    \item 证明 $(1,1)$ 是 $A\times B$ 的单位元；
    \item 证明 $(a,b)$ 的逆元是 $(a^{-1},b^{-1})$。
  \end{enumerate}
  \item 证明 $A\times B$ 是阿贝尔群，当且仅当 $A$ 与 $B$ 都是阿贝尔群。
  \item 证明 $A\times B$ 中的元素 $(a,1)$ 与 $(1,b)$ 交换，并推出 $(a,b)$ 的阶是 $|a|$ 与 $|b|$ 的最小公倍数。
  \item 证明任一偶数阶有限群 $G$ 含有阶为 $2$ 的元素。［令 $t(G)=\{g\in G\mid g\neq g^{-1}\}$。证明 $t(G)$ 的元素个数为偶数，并且 $G-t(G)$ 中每个非单位元素的阶为 $2$。］
  \item 若 $x$ 是 $G$ 中有限阶 $n$ 的元素，证明 $1,x,x^2,\ldots,x^{n-1}$ 两两不同。推出 $|x|\leq |G|$。
  \item 设 $x$ 是 $G$ 中有限阶 $n$ 的元素。
  \begin{enumerate}[label=(\alph*)]
    \item 证明若 $n$ 为奇数，则对所有 $i=1,2,\ldots,n-1$，有 $x^i\neq x^{-i}$。
    \item 证明若 $n=2k$ 且 $1\leq i<n$，则 $x^i=x^{-i}$ 当且仅当 $i=k$。
  \end{enumerate}
  \item 若 $x$ 是 $G$ 中无限阶元素，证明元素 $x^n$（$n\in\ZZ$）两两不同。
  \item 若 $x$ 是 $G$ 中有限阶 $n$ 的元素，用带余除法证明 $x$ 的任意整数幂都等于集合 $\{1,x,x^2,\ldots,x^{n-1}\}$ 中的某个元素（因此这些就是由 $x$ 生成的循环子群的全部不同元素，参见上面的练习 27）。
  \item 假设 $G=\{1,a,b,c\}$ 是阶为 $4$、单位元为 $1$ 的群。再假设 $G$ 没有阶为 $4$ 的元素（所以由练习 32，每个元素阶都 $\leq 3$）。用消去律证明 $G$ 的群表唯一。推出 $G$ 是阿贝尔群。
\end{enumerate}

\section{二面体群}

群的一类重要例子是几何对象的对称所组成的群。最简单的子类是几何对象为正平面图形的情形。

对每个 $n\in\ZZ^+$，$n\geq3$，令 $D_{2n}$ 为正 $n$ 边形的所有对称组成的集合。这里的对称是指正 $n$ 边形的任一刚体运动，它可以通过取该正 $n$ 边形的一个副本，在三维空间中以任意方式移动它，然后把该副本放回原正 $n$ 边形上，使之恰好覆盖原图形来实现。更精确地说，我们可以先给 $n$ 个顶点编号，例如如下图所示，由此描述这些对称。

\begin{center}
\begin{tikzpicture}[scale=1.0]
  \def\n{7}
  \foreach \i in {1,...,7} {
    \coordinate (P\i) at ({90-360*(\i-1)/7}:1.6);
    \fill (P\i) circle (1.2pt);
    \node at ({90-360*(\i-1)/7}:1.95) {$\i$};
  }
  \draw (P1)--(P2)--(P3)--(P4)--(P5)--(P6)--(P7)--cycle;
\end{tikzpicture}
\end{center}

于是每个对称 $s$ 都可由 $\{1,2,3,\ldots,n\}$ 的相应置换 $\sigma$ 唯一描述：若对称 $s$ 把顶点 $i$ 放到原来顶点 $j$ 所在的位置，则 $\sigma$ 是把 $i$ 送到 $j$ 的置换。例如，若 $s$ 是绕 $n$ 边形中心顺时针旋转 $2\pi/n$ 弧度，则 $\sigma$ 把 $i$ 送到 $i+1$，$1\leq i\leq n-1$，并且 $\sigma(n)=1$。现在把 $D_{2n}$ 作成一个群：对 $s,t\in D_{2n}$，定义 $st$ 为先对 $n$ 边形施行 $t$，再施行 $s$ 所得到的对称（注意我们把对称看作 $n$ 边形上的函数，所以 $st$ 就是函数复合——按通常方式从右向左读）。若 $s,t$ 在顶点上分别实现置换 $\sigma,\tau$，则 $st$ 实现 $\sigma\circ\tau$。$D_{2n}$ 上的二元运算是结合的，因为函数复合是结合的。$D_{2n}$ 的单位元是恒等对称（它固定所有顶点），记作 $1$；$s\in D_{2n}$ 的逆元是把 $s$ 的所有刚体运动反向的对称（所以若 $s$ 在顶点上实现置换 $\sigma$，则 $s^{-1}$ 实现 $\sigma^{-1}$）。下一段我们说明
\[
  |D_{2n}|=2n,
\]
因此 $D_{2n}$ 称为\emph{阶为 $2n$ 的二面体群}。有些教材把这个群记为 $D_n$；不过在群论文献中，$D_{2n}$（下标给出群的阶而不是顶点数）更常见。

为求阶 $|D_{2n}|$，注意给定任一顶点 $i$，存在一个对称把顶点 $1$ 送到位置 $i$。由于顶点 $2$ 与顶点 $1$ 相邻，顶点 $2$ 必须落在位置 $i+1$ 或 $i-1$（其中 $n+1$ 读作 $1$，$1-1$ 读作 $n$，即顶点编号按模 $n$ 读取）。而且，在第一个对称之后再复合关于通过顶点 $i$ 和 $n$ 边形中心的直线的反射，可见顶点 $2$ 可由某个对称送到位置 $i+1$ 或 $i-1$ 中任意一个。因此，应用对称后，有序顶点对 $1,2$ 有 $n\cdot2$ 种可能位置。由于对称是刚体运动，一旦有序顶点对 $1,2$ 的位置被指定，该对称在所有其余顶点上的作用就完全确定。因此正 $n$ 边形恰有 $2n$ 个对称。而且我们可以显式列出这 $2n$ 个对称：绕中心旋转 $2\pi i/n$ 弧度的 $n$ 个旋转，$0\leq i\leq n-1$，以及关于 $n$ 条对称轴的 $n$ 个反射（若 $n$ 为奇数，每条对称轴通过一个顶点和对边的中点；若 $n$ 为偶数，则有 $n/2$ 条对称轴通过两个相对顶点，有 $n/2$ 条对称轴垂直平分两条相对边）。例如，若 $n=4$，并在 $x,y$ 平面中把正方形画在原点处，则对称轴为
\[
  x=0\quad(y\text{-轴}),\qquad y=0\quad(x\text{-轴}),\qquad y=x,
  \qquad y=-x,
\]
（注意，关于原点的“反射”并不是反射，而是旋转 $\pi$ 弧度）。

\begin{center}
\begin{tikzpicture}[scale=1.2]
  \draw[->] (-1.6,0)--(1.7,0) node[right] {$x$};
  \draw[->] (0,-1.6)--(0,1.7) node[above] {$y$};
  \draw (-1,1)--(1,1)--(1,-1)--(-1,-1)--cycle;
  \node at (1.15,1.15) {$1$};
  \node at (1.15,-1.15) {$2$};
  \node at (-1.15,-1.15) {$3$};
  \node at (-1.15,1.15) {$4$};
  \draw[dashed] (-1.4,-1.4)--(1.4,1.4) node[right] {$y=x$};
  \draw[dashed] (-1.4,1.4)--(1.4,-1.4) node[right] {$y=-x$};
\end{tikzpicture}
\end{center}

由于二面体群将在全文中被大量用作例子，我们固定一些记号，并提及一些计算；它们将简化以后的计算，也有助于把 $D_{2n}$ 看作一个抽象群（而不必每次都回到几何背景）。固定一个以原点为中心、位于 $x,y$ 平面中的正 $n$ 边形，并按顺时针方向依次给顶点标号 $1$ 到 $n$。令 $r$ 为绕原点顺时针旋转 $2\pi/n$ 弧度的旋转。令 $s$ 为关于通过顶点 $1$ 和原点的对称轴的反射（对每个 $n$ 都用同样的字母，但上下文总会使 $n$ 明确）。下面计算的细节留作练习（我们大多将处理 $D_6$ 和 $D_8$，读者也许可以先对 $n=3$ 和 $n=4$ 试做这些练习）：
\begin{enumerate}[label=(\arabic*)]
  \item $1,r,r^2,\ldots,r^{n-1}$ 两两不同，并且 $r^n=1$，所以 $|r|=n$。
  \item $|s|=2$。
  \item 对任意 $i$，$s\neq r^i$。
  \item 若 $0\leq i,j\leq n-1$ 且 $i\neq j$，则 $sr^i\neq sr^j$，所以
  \[
    D_{2n}=\{1,r,r^2,\ldots,r^{n-1},s,sr,sr^2,\ldots,sr^{n-1}\},
  \]
  即每个元素可唯一写成 $s^kr^i$ 的形式，其中 $k=0$ 或 $1$，且 $0\leq i\leq n-1$。
  \item $rs=sr^{-1}$。［先求出 $s$ 在 $\{1,2,\ldots,n\}$ 上实现的置换，再分别计算此等式两边对顶点 $1$ 和 $2$ 的作用。］这尤其说明 $r$ 与 $s$ 不交换，因此 $D_{2n}$ 是非阿贝尔的。
  \item 对所有 $0\leq i\leq n$，有 $r^is=sr^{-i}$。［对 $i$ 作归纳，并利用 $r^{i+1}s=r(r^is)$ 以及前面的计算。］这说明如何把 $s$ 与 $r$ 的幂交换位置。
\end{enumerate}

完成这些计算后，我们现在观察到，$D_{2n}$ 的完整乘法表可以仅用 $r$ 和 $s$ 写出；也就是说，$D_{2n}$ 的所有元素都可唯一表示为 $s^kr^i$，其中 $k=0$ 或 $1$，$0\leq i\leq n-1$，而且任意两个这种形式的元素的乘积，都可以只用“关系” (1)、(2)、(6)（把所有指数模 $n$ 化简）化为另一个这种形式的元素。例如，若 $n=12$，则
\[
  (sr^9)(sr^6)=s(r^9s)r^6=s(sr^{-9})r^6=s^2r^{-9+6}=r^{-3}=r^9.
\]

\subsection*{生成元与关系}

二面体群中生成元 $r$ 与 $s$ 的使用，为在 $D_{2n}$ 中计算提供了一种简单而简洁的方式。我们可以类似地为任意群引入生成元与关系的概念。较早拥有这些概念（在其形式化证明之前）是有用的，因为它们提供了描述和计算许多群的简单方式。生成元将在第 2.4 节中更详细地讨论；当第 6.3 节引入自由群概念时，这两个概念都会得到严格处理。

群 $G$ 的元素子集 $S$ 若具有如下性质：$G$ 的每个元素都可以写成 $S$ 中元素及其逆元的一个（有限）乘积，则称 $S$ 是 $G$ 的一个\emph{生成元集}。记号上我们写作 $G=\langle S\rangle$，并说 $G$ \emph{由 $S$ 生成}，或 $S$ \emph{生成} $G$。例如，整数 $1$ 是加法群 $\ZZ$ 的一个生成元，因为每个整数都是有限多个 $+1$ 与 $-1$ 的和，所以 $\ZZ=\langle 1\rangle$。由 $D_{2n}$ 的性质 (4)，集合 $S=\{r,s\}$ 是 $D_{2n}$ 的一个生成元集，所以 $D_{2n}=\langle r,s\rangle$。我们稍后将看到，在有限群 $G$ 中，若 $G$ 的每个元素都是 $S$ 中元素的有限乘积，则 $S$ 生成 $G$（即不必另外把 $S$ 中元素的逆元也包括进来）。

在一般群 $G$ 中，生成元满足的任何方程称为 $G$ 中的\emph{关系}。因此在 $D_{2n}$ 中有关系：$r^n=1$、$s^2=1$ 和 $rs=sr^{-1}$。此外，在 $D_{2n}$ 中这三个关系还具有额外性质：群元素之间的任何其他关系都可由这三个关系推出（这并非立即显然；它来自这样一个事实：只用这三个关系就能确定两个群元素何时相等）。

一般地，若某个群 $G$ 由子集 $S$ 生成，并且有一组关系，例如 $R_1,R_2,\ldots,R_m$（这里每个 $R_i$ 都是由 $S\cup\{1\}$ 中元素写成的一个方程），使得 $S$ 中元素之间的任一关系都能由这些关系推出，则称这些生成元与关系给出了 $G$ 的一个\emph{表示}，并写作
\[
  G=\langle S\mid R_1,R_2,\ldots,R_m\rangle.
\]
因此，二面体群 $D_{2n}$ 的一个表示（使用上述生成元和关系）是
\begin{equation}
  D_{2n}=\langle r,s\mid r^n=s^2=1,\ rs=sr^{-1}\rangle. \tag{1.1}
\end{equation}
我们将看到，用这个表示描述 $D_{2n}$（而不是总回到原始的几何描述）会大大简化对这些群的处理。

表示给出了描述许多群的一种容易方法，但也有若干微妙之处需要考虑。其中之一是：在任意表示中，可能很难（甚至不可能）判断两个用给定生成元表达的群元素是否相等。因此，所表示的群的阶是多少，甚至该群是有限还是无限，可能都不明显。例如，可以证明
\[
  \langle x_1,y_1\mid x_1^2=y_1^2=(x_1y_1)^2=1\rangle
\]
是一个阶为 $4$ 的群的表示，而
\[
  \langle x_2,y_2\mid x_2^3=y_2^3=(x_2y_2)^3=1\rangle
\]
则是一个无限群的表示（参见练习）。

另一个微妙之处是，即使在相当简单的表示中，也可能因为关系以某种不明显的方式交织在一起而发生某些“塌缩”；也就是说，可能存在没有在表示中显式给出、而是指定关系的后果的“隐藏”或隐含关系。一般而言，这种塌缩使得甚至很难确定所表示群大小的下界。例如，假设模仿 $D_{2n}$ 的表示来试图创造另一个群，定义
\begin{equation}
  X_{2n}=\langle x,y\mid x^n=y^2=1,\ xy=yx^2\rangle. \tag{1.2}
\end{equation}
“交换”关系 $xy=yx^2$ 决定了如何交换 $y$ 和 $x$（即如何把 $y$ 从 $x$ 的右侧“移到”左侧），所以与 $D_{2n}$ 中一样，这个群的每个元素都可写成 $y^kx^i$ 的形式，所有 $y$ 的幂在左，所有 $x$ 的幂在右。此外，由前两个关系，$x$ 和 $y$ 的任意幂都可化简，使 $i$ 位于 $0$ 与 $n-1$ 之间，并且 $k$ 为 $0$ 或 $1$。因此人们也许会猜测 $X_{2n}$ 仍是一个阶为 $2n$ 的群。但事实并非如此，因为在这个群中存在一个“隐藏”关系，它由关系 $x=xy^2$（因为 $y^2=1$）通过反复应用交换关系和结合律把 $y$ 移到左侧而得到：
\begin{align*}
  x&=x1=(xy)y=(yx^2)y=(yx)(xy)=(yx)(yx^2)\
   &=y(xy)x^2=y(yx^2)x^2=1x^4=x^4.
\end{align*}
由于 $x^4=x$，由消去律得 $x^3=1$ 于 $X_{2n}$ 中；由上面的讨论可知，对任意 $n$，$X_{2n}$ 的阶至多为 $6$。根据 $n$ 的取值，还可能发生更多塌缩（见练习）。

再例如，考虑表示
\begin{equation}
  Y=\langle u,v\mid u^4=v^3=1,\ uv=v^2u^2\rangle. \tag{1.3}
\end{equation}
在这个情形中，容易猜测 $Y$ 是阶为 $12$ 的群，但仍有额外的隐含关系。事实上，这个群 $Y$ 退化为阶为 $1$ 的平凡群；也就是说，$u$ 和 $v$ 满足额外关系 $u=1$ 与 $v=1$（证明概要见练习）。

这种塌缩不会发生在 $D_{2n}$ 的表示中，因为我们已经用独立的（几何）方法说明存在一个阶为 $2n$、由 $r$ 与 $s$ 生成并满足 (1) 中关系的群。因此，只具有这些关系的群的阶至少为 $2n$。另一方面，容易看出（使用上面对 $X_{2n}$ 的同类论证以及交换关系 $rs=sr^{-1}$），任意由 (1) 中的生成元与关系定义的群，其阶至多为 $2n$。于是，具有表示 (1) 的群阶恰为 $2n$，并且这个群确实就是正 $n$ 边形的对称群。

我们对表示 (1) 所拥有的额外信息，是存在一个已知阶的群满足这些信息。相比之下，对满足 (2) 或 (3) 中关系的任何群，我们没有独立的知识。没有这种独立的“下界”信息，我们甚至可能无法判断给定表示是否只是描述了平凡群，如 (3) 所示。

虽然一般来说，用表示来规定群时必须极其小心，但对已知群使用表示是一种强有力的概念工具和计算工具。关于表示的更多结果，包括更复杂的例子，见第 6.3 节。

\subsection*{练习}

在这些练习中，$D_{2n}$ 取通常的表示
\[
  D_{2n}=\langle r,s\mid r^n=s^2=1,\ rs=sr^{-1}\rangle.
\]
\begin{enumerate}[label=\arabic*.]
  \item 计算下列群中每个元素的阶：
  \begin{enumerate}[label=(\alph*)]
    \item $D_6$；\quad (b) $D_8$；\quad (c) $D_{10}$。
  \end{enumerate}
  \item 用上面的生成元与关系证明：若 $x$ 是 $D_{2n}$ 中任一不是 $r$ 的幂的元素，则 $rx=xr^{-1}$。
  \item 用上面的生成元与关系证明：$D_{2n}$ 中任一不是 $r$ 的幂的元素都有阶 $2$。推出 $D_{2n}$ 由两个元素 $s$ 与 $sr$ 生成，且这两个元素均有阶 $2$。
  \item 若 $n=2k$ 为偶数且 $n\geq4$，证明 $z=r^k$ 是一个阶为 $2$ 且与 $D_{2n}$ 的所有元素交换的元素。还要证明 $z$ 是 $D_{2n}$ 中唯一与所有元素交换的非单位元素。［参见第 1 节练习 33。］
  \item 若 $n$ 为奇数且 $n\geq3$，证明单位元是 $D_{2n}$ 中唯一与所有元素交换的元素。［参见第 1 节练习 33。］
  \item 设 $x$ 与 $y$ 是任意群 $G$ 中两个阶为 $2$ 的元素。证明若 $t=xy$，则 $tx=xt^{-1}$（所以若 $n=|xy|<\infty$，则 $x,t$ 在 $G$ 中满足与 $D_{2n}$ 中 $s,r$ 相同的关系）。
  \item 证明
  \[
    \langle a,b\mid a^2=b^2=(ab)^n=1\rangle
  \]
  给出了 $D_{2n}$ 关于练习 3 中得到的两个阶为 $2$ 的生成元 $a=s$ 与 $b=sr$ 的表示。［证明 $r,s$ 的关系可由 $a,b$ 的关系推出；反过来，$a,b$ 的关系也可由 $r,s$ 的关系推出。］
  \item 求 $D_{2n}$ 中由 $r$ 生成的循环子群的阶（参见第 1 节练习 27）。
\end{enumerate}

在练习 9 到 13 中，你可以仿照求 $D_{2n}$ 阶的证明来求给定柏拉图立体在 $\RR^3$ 中的刚体运动群（也称旋转群）的阶：求一对相邻顶点可被送到的位置数。或者，也可以先求一个给定面可被送到的位置数，再在固定一个面后求该面上一顶点可被送到的位置数。
\begin{enumerate}[label=\arabic*.,resume]
  \item 设 $G$ 是正四面体在 $\RR^3$ 中的刚体运动群。证明 $|G|=12$。
  \item 设 $G$ 是立方体在 $\RR^3$ 中的刚体运动群。证明 $|G|=24$。
  \item 设 $G$ 是正八面体在 $\RR^3$ 中的刚体运动群。证明 $|G|=24$。
  \item 设 $G$ 是正十二面体在 $\RR^3$ 中的刚体运动群。证明 $|G|=60$。
  \item 设 $G$ 是正二十面体在 $\RR^3$ 中的刚体运动群。证明 $|G|=60$。
  \item 找出 $\ZZ$ 的一个生成元集。
  \item 找出 $\ZZ/n\ZZ$ 的一组生成元与关系。
  \item 证明群 $\langle x_1,y_1\mid x_1^2=y_1^2=(x_1y_1)^2=1\rangle$ 是二面体群 $D_4$（其中可把 $x_1$ 改写为字母 $r$，把 $y_1$ 改写为 $s$）。［证明最后一个关系等价于 $x_1y_1=y_1x_1^{-1}$。］
  \item 设 $X_{2n}$ 是由 (1.2) 中表示给出的群。
  \begin{enumerate}[label=(\alph*)]
    \item 证明若 $n=3k$，则 $X_{2n}$ 的阶为 $6$；当把 $x$ 换成 $r$、$y$ 换成 $s$ 时，它与 $D_6$ 有相同的生成元与关系。
    \item 证明若 $(3,n)=1$，则 $x$ 满足额外关系 $x=1$。在此情形下推出 $X_{2n}$ 的阶为 $2$。［用 $x^n=1$ 与 $x^3=1$。］
  \end{enumerate}
  \item 设 $Y$ 是由 (1.3) 中表示给出的群。
  \begin{enumerate}[label=(\alph*)]
    \item 证明 $v^2=v^{-1}$。［用关系 $v^3=1$。］
    \item 证明 $v$ 与 $u^3$ 交换。［把左边写为 $(v^2u^2)(uv)$ 并用关系把它化为右边，从而证明 $v^2u^3v=u^3$；再用 (a)。］
    \item 证明 $v$ 与 $u$ 交换。［证明 $u^9=u$，然后用 (b)。］
    \item 证明 $uv=1$。［用 (c) 与最后一个关系。］
    \item 证明 $u=1$，推出 $v=1$，并得出 $Y=1$。［用 (d) 与等式 $u^4v^3=1$。］
  \end{enumerate}
\end{enumerate}


\section{对称群}

令 $\Omega$ 为任一非空集合，令 $S_\Omega$ 为从 $\Omega$ 到自身的所有双射组成的集合（即 $\Omega$ 的所有置换组成的集合）。集合 $S_\Omega$ 在函数复合 $\circ$ 下是群。注意，$\circ$ 是 $S_\Omega$ 上的二元运算，因为若 $\sigma:\Omega\to\Omega$ 和 $\tau:\Omega\to\Omega$ 都是双射，则 $\sigma\circ\tau$ 也是从 $\Omega$ 到 $\Omega$ 的双射。由于函数复合一般是结合的，$\circ$ 是结合的。$S_\Omega$ 的单位元是置换 $1$，定义为对所有 $a\in\Omega$，$1(a)=a$。每个置换 $\sigma$ 都有一个（双侧）逆函数 $\sigma^{-1}:\Omega\to\Omega$，满足 $\sigma\circ\sigma^{-1}=\sigma^{-1}\circ\sigma=1$。因此，$(S_\Omega,\circ)$ 满足所有群公理。这个群称为集合 $\Omega$ 上的\emph{对称群}。重要的是要认识到，$S_\Omega$ 的元素是 $\Omega$ 的置换，而不是 $\Omega$ 本身的元素。

在特殊情形 $\Omega=\{1,2,3,\ldots,n\}$ 下，$\Omega$ 上的对称群记为 $S_n$，称为\emph{$n$ 次对称群}。\footnote{我们将在第 6 节看到，$S_\Omega$ 的结构只依赖于 $\Omega$ 的基数，而不依赖于 $\Omega$ 的具体元素；因此若 $\Omega$ 是任一含 $n$ 个元素的有限集，则 $S_\Omega$ “看起来像” $S_n$。}群 $S_n$ 将在全文中发挥重要作用：它本身是一个相当重要的研究对象，同时也是说明和激发一般理论的工具。

首先说明 $S_n$ 的阶为 $n!$。集合 $\{1,2,3,\ldots,n\}$ 的置换恰好就是该集合到自身的单射，因为该集合有限（命题 0.1）；于是我们可以数出单射的个数。单射 $\sigma$ 可把数字 $1$ 送到 $\{1,2,3,\ldots,n\}$ 中任意一个元素；$\sigma(2)$ 随后可为该集合中除 $\sigma(1)$ 外的任一元素（所以 $\sigma(2)$ 有 $n-1$ 种选择）；$\sigma(3)$ 可为除 $\sigma(1)$ 和 $\sigma(2)$ 外的任一元素（所以有 $n-2$ 种选择）；依此类推。因此，从 $\{1,2,3,\ldots,n\}$ 到自身的单射恰有
\[
  n(n-1)(n-2)\cdots 2\cdot1=n!
\]
个。因此 $\{1,2,3,\ldots,n\}$ 恰有 $n!$ 个置换，也就是 $S_n$ 中恰有 $n!$ 个元素。

现在描述一种写出 $S_n$ 的元素 $\sigma$ 的有效记号，全文都将使用它，这称为\emph{循环分解}。

一个\emph{循环}是一串整数，用来表示 $S_n$ 中把这些整数循环置换（并固定所有其他整数）的元素。循环 $(a_1\ a_2\ \cdots\ a_m)$ 是这样一个置换：把 $a_i$ 送到 $a_{i+1}$，$1\leq i\leq m-1$，并把 $a_m$ 送到 $a_1$。例如，$(2\ 1\ 3)$ 是把 $2$ 映到 $1$、把 $1$ 映到 $3$、把 $3$ 映到 $2$ 的置换。一般地，对每个 $\sigma\in S_n$，数字 $1$ 到 $n$ 会被重新排列并分组成 $k$ 个循环，形如
\[
  (a_1\ a_2\ \cdots\ a_{m_1})
  (a_{m_1+1}\ \cdots\ a_{m_2})\cdots
  (a_{m_{k-1}+1}\ \cdots\ a_{m_k}),
\]
其中 $m_k=n$；从这个表达式可以很容易读出 $\sigma$ 对 $1$ 到 $n$ 中任一数字的作用。对任意 $x\in\{1,2,3,\ldots,n\}$，先在上述表达式中找到 $x$。若 $x$ 后面不是紧跟一个右括号（即 $x$ 不是某个循环的右端），则 $\sigma(x)$ 是紧挨着 $x$ 右侧的整数。若 $x$ 后面紧跟右括号，则 $\sigma(x)$ 是以 $x$ 结尾的那个循环开头的数（即若 $x=a_{m_i}$，则 $\sigma(x)=a_{m_{i-1}+1}$，其中 $m_0=0$）。所有循环的乘积称为 $\sigma$ 的\emph{循环分解}。

下面给出求 $S_n$ 中元素 $\sigma$ 的循环分解的算法，并通过一个具体置换演示该算法。算法的证明以及循环分解唯一性的完整分析推迟到第 4 章。

令 $n=13$，定义 $\sigma\in S_{13}$ 为
\begin{align*}
&\sigma(1)=12,\quad \sigma(2)=13,\quad \sigma(3)=3,\quad \sigma(4)=1,\quad \sigma(5)=11,\\
&\sigma(6)=9,\quad \sigma(7)=5,\quad \sigma(8)=10,\quad \sigma(9)=6,\quad \sigma(10)=4,\\
&\sigma(11)=7,\quad \sigma(12)=8,\quad \sigma(13)=2.
\end{align*}

\paragraph{循环分解算法}
\begin{center}
\begin{tabular}{p{0.56\linewidth}|p{0.34\linewidth}}
\textbf{方法} & \textbf{例子} \\ \hline
要开始一个新循环，选取 $\{1,2,\ldots,n\}$ 中尚未在前面循环中出现的最小元素，称之为 $a$（若刚开始，则 $a=1$）；开始新的循环：$(a$
& $(1$ \\[0.8em]
从 $\sigma$ 的给定描述中读出 $\sigma(a)$，称之为 $b$。若 $b=a$，则用右括号结束该循环（不写下 $b$）；这完成一个循环，回到第 1 步。若 $b\neq a$，则把 $b$ 写在该循环中 $a$ 的旁边：$(a\ b$
& $\sigma(1)=12=b$，$12\neq1$，所以写成：$(1\ 12$ \\[0.8em]
从 $\sigma$ 的给定描述中读出 $\sigma(b)$，称之为 $c$。若 $c=a$，则用右括号结束该循环，从而完成该循环，并回到第 1 步。若 $c\neq a$，则把 $c$ 写在该循环中 $b$ 的旁边：$(a\ b\ c$。以数字 $c$ 作为新的 $b$，重复这一步，直到循环闭合。
& $\sigma(12)=8$，$8\neq1$，所以继续该循环为：$(1\ 12\ 8$ 
\end{tabular}
\end{center}
自然地，当 $\{1,2,\ldots,n\}$ 中所有数都已出现在某个循环中时，这个过程停止。对该例中的特定 $\sigma$，这给出
\[
  \sigma=(1\ 12\ 8\ 10\ 4)(2\ 13)(3)(5\ 11\ 7)(6\ 9).
\]
一个循环的\emph{长度}是出现在其中的整数个数。长度为 $t$ 的循环称为一个 $t$-循环。若两个循环没有公共数字，则称它们\emph{不交}。因此，上面的元素 $\sigma$ 是 $5$ 个两两不交循环的乘积：一个 $5$-循环、一个 $2$-循环、一个 $1$-循环、一个 $3$-循环和另一个 $2$-循环。

今后我们约定不写出 $1$-循环。因此，若某个整数 $i$ 没有出现在置换 $\tau$ 的循环分解中，则默认 $\tau(i)=i$，即 $\tau$ 固定 $i$。$S_n$ 的恒等置换的循环分解是 $(1)(2)\cdots(n)$，并将简写为 $1$。所以算法的最后一步是：
\[
\boxed{\text{最后一步：删去所有长度为 }1\text{ 的循环}}
\]
所以该例中 $\sigma$ 的循环分解最终写为
\[
  \sigma=(1\ 12\ 8\ 10\ 4)(2\ 13)(5\ 11\ 7)(6\ 9).
\]
这个约定的优点是，$S_n$ 中元素 $\tau$ 的循环分解，同时也是 $S_m$（$m\geq n$）中如下置换的循环分解：它在 $\{1,2,3,\ldots,n\}$ 上按 $\tau$ 作用，并固定 $\{n+1,n+2,\ldots,m\}$ 中每个元素。因此，例如 $(1\ 2)$ 表示交换 $1$ 和 $2$ 并固定所有更大整数的置换，无论把它看作 $S_2$、$S_3$、$S_4$ 等中的元素。

再例如，$S_3$ 的 $6$ 个元素具有如下循环分解：
\[
\begin{array}{c|c|c}
 & \text{置换的取值} & \text{循环分解}\\ \hline
\sigma_1&1\mapsto1,\ 2\mapsto2,\ 3\mapsto3&1\\
\sigma_2&1\mapsto1,\ 2\mapsto3,\ 3\mapsto2&(2\ 3)\\
\sigma_3&1\mapsto3,\ 2\mapsto2,\ 3\mapsto1&(1\ 3)\\
\sigma_4&1\mapsto2,\ 2\mapsto1,\ 3\mapsto3&(1\ 2)\\
\sigma_5&1\mapsto2,\ 2\mapsto3,\ 3\mapsto1&(1\ 2\ 3)\\
\sigma_6&1\mapsto3,\ 2\mapsto1,\ 3\mapsto2&(1\ 3\ 2)
\end{array}
\]
对任意 $\sigma\in S_n$，$\sigma^{-1}$ 的循环分解可通过把 $\sigma$ 的循环分解中每个循环里的数字按相反顺序写出得到。例如，若 $\sigma=(1\ 12\ 8\ 10\ 4)(2\ 13)(5\ 11\ 7)(6\ 9)$，则
\[
  \sigma^{-1}=(4\ 10\ 8\ 12\ 1)(13\ 2)(7\ 11\ 5)(9\ 6).
\]

在 $S_n$ 中计算乘积是直接的，只要记住在计算 $\sigma\circ\tau$ 时从右向左读置换。只需“跟踪”元素在连续置换下的去向。例如，在乘积
\[
  (1\ 2\ 3)\circ(1\ 2)(3\ 4)
\]
中，数字 $1$ 先被第一个置换送到 $2$，然后 $2$ 被第二个置换送到 $3$，所以复合映射把 $1$ 送到 $3$。为求乘积的循环分解，接下来要看 $3$ 的去向：$3$ 先被送到 $4$，而后 $4$ 被固定，所以 $3$ 被复合映射送到 $4$。类似地，$4$ 先被送到 $3$，再被送到 $1$，于是乘积中的这个循环闭合为 $(1\ 3\ 4)$。最后，$2$ 被送到 $1$，再被送到 $2$，所以 $2$ 被这个乘积固定。因此
\[
  (1\ 2\ 3)\circ(1\ 2)(3\ 4)=(1\ 3\ 4).
\]
作为更多例子，
\[
  (1\ 2)\circ(1\ 3)=(1\ 3\ 2),\qquad (1\ 3)\circ(1\ 2)=(1\ 2\ 3).
\]
这特别说明，对所有 $n\geq3$，$S_n$ 都是非阿贝尔群。

循环分解中的每个循环 $(a_1\ a_2\ \cdots\ a_m)$ 可看成把 $a_1,a_2,\ldots,a_m$ 循环置换并固定所有其他整数的置换。由于不交循环置换的数字属于彼此不交的集合，所以
\[
  \text{不交循环彼此交换。}
\]
因此，任意不交循环乘积（特别是循环分解）中循环的排列顺序不会改变该置换。

此外，由于一个给定循环 $(a_1\ a_2\ \cdots\ a_m)$ 循环置换集合 $\{a_1,a_2,\ldots,a_m\}$，所以循环中数字本身可以循环移位而不改变置换，即
\[
  (a_1\ a_2\ \cdots\ a_m)=(a_2\ a_3\ \cdots\ a_m\ a_1)=(a_3\ a_4\ \cdots\ a_m\ a_1\ a_2)=\cdots=(a_m\ a_1\ a_2\ \cdots\ a_{m-1}).
\]
因此，例如 $(1\ 2)=(2\ 1)$，且 $(1\ 2\ 3\ 4)=(3\ 4\ 1\ 2)$。按惯例，通常把循环中出现的最小数字写在最前面。

使用循环时必须稍加小心，因为一个置换可用许多方式写成循环的任意乘积。例如在 $S_3$ 中，
\[
  (1\ 2\ 3)=(1\ 2)(2\ 3)=(1\ 3)(1\ 3\ 2)(1\ 3)
\]
等等。但是（我们将会证明）每个置换的循环分解是把它表示成不交循环之积的唯一方式（只差循环的重新排序和每个循环内部数字的循环移位）。把任意循环乘积化为不交循环之积，使我们能够一眼判断两个置换是否相同。这种记号的另一个优点是：证明一个置换的阶等于其循环分解中各循环长度的最小公倍数，是一个练习（下面给出提示）。

\subsection*{练习}
\begin{enumerate}[label=\arabic*.]
  \item 设 $\sigma$ 是置换
  \[
  1\mapsto3,\\ 2\mapsto4,\\ 3\mapsto5,\\ 4\mapsto2,\\ 5\mapsto1,
  \]
  并设 $\tau$ 是置换
  \[
  1\mapsto5,\\ 2\mapsto3,\\ 3\mapsto2,\\ 4\mapsto4,\\ 5\mapsto1.
  \]
  求下列置换的循环分解：$\sigma,\tau,\sigma^2,\sigma\tau,\tau\sigma,\tau^2\sigma$。
  \item 设 $\sigma$ 是置换
  \[
  \begin{array}{ccccc}
  1\mapsto13&2\mapsto2&3\mapsto15&4\mapsto14&5\mapsto10\\
  6\mapsto6&7\mapsto12&8\mapsto3&9\mapsto4&10\mapsto1\\
  11\mapsto7&12\mapsto9&13\mapsto5&14\mapsto11&15\mapsto8
  \end{array}
  \]
  并设 $\tau$ 是置换
  \[
  \begin{array}{ccccc}
  1\mapsto14&2\mapsto9&3\mapsto10&4\mapsto2&5\mapsto12\\
  6\mapsto6&7\mapsto5&8\mapsto11&9\mapsto15&10\mapsto3\\
  11\mapsto8&12\mapsto7&13\mapsto4&14\mapsto1&15\mapsto13.
  \end{array}
  \]
  求下列置换的循环分解：$\sigma,\tau,\sigma^2,\sigma\tau,\tau\sigma,\tau^2\sigma$。
  \item 对前两题中求得循环分解的每个置换，计算其阶。
  \item 计算下列群中每个元素的阶：\begin{enumerate}[label=(\alph*)]
    \item $S_3$；\quad \item $S_4$。
  \end{enumerate}
  \item 求 $(1\ 12\ 8\ 10\ 4)(2\ 13)(5\ 11\ 7)(6\ 9)$ 的阶。
  \item 写出 $S_4$ 中每个阶为 $4$ 的元素的循环分解。
  \item 写出 $S_4$ 中每个阶为 $2$ 的元素的循环分解。
  \item 证明若 $\Omega=\{1,2,3,\ldots\}$，则 $S_\Omega$ 是无限群（不要说 $\infty!=\infty$）。
  \item
  \begin{enumerate}[label=(\alph*)]
    \item 令 $\sigma$ 为 $12$-循环 $(1\ 2\ 3\ 4\ 5\ 6\ 7\ 8\ 9\ 10\ 11\ 12)$。对哪些正整数 $i$，$\sigma^i$ 仍是 $12$-循环？
    \item 令 $\tau$ 为 $8$-循环 $(1\ 2\ 3\ 4\ 5\ 6\ 7\ 8)$。对哪些正整数 $i$，$\tau^i$ 仍是 $8$-循环？
    \item 令 $\omega$ 为 $14$-循环 $(1\ 2\ 3\ 4\ 5\ 6\ 7\ 8\ 9\ 10\ 11\ 12\ 13\ 14)$。对哪些正整数 $i$，$\omega^i$ 仍是 $14$-循环？
  \end{enumerate}
  \item 证明若 $\sigma$ 是 $m$-循环 $(a_1\ a_2\ \cdots\ a_m)$，则对所有 $i\in\{1,2,\ldots,m\}$，$\sigma^i(a_k)=a_{k+i}$，其中当 $k+i>m$ 时，把 $k+i$ 换成模 $m$ 的最小剩余。推出 $|\sigma|=m$。
  \item 令 $\sigma$ 为 $m$-循环 $(1\ 2\ \cdots\ m)$。证明 $\sigma^i$ 也是 $m$-循环，当且仅当 $i$ 与 $m$ 互素。
  \item
  \begin{enumerate}[label=(\alph*)]
    \item 若 $\tau=(1\ 2)(3\ 4)(5\ 6)(7\ 8)(9\ 10)$，判断是否存在某个 $n$-循环 $\sigma$（$n\geq10$）以及某个整数 $k$，使得 $\tau=\sigma^k$。
    \item 若 $\tau=(1\ 2)(3\ 4\ 5)$，判断是否存在某个 $n$-循环 $\sigma$（$n\geq5$）以及某个整数 $k$，使得 $\tau=\sigma^k$。
  \end{enumerate}
  \item 证明 $S_n$ 中一个元素阶为 $2$，当且仅当它的循环分解是一些可交换的 $2$-循环的乘积。
  \item 设 $p$ 是素数。证明 $S_n$ 中一个元素阶为 $p$，当且仅当它的循环分解是一些可交换的 $p$-循环的乘积。用一个显式例子说明若 $p$ 不是素数，则这不一定成立。
  \item 证明 $S_n$ 中元素的阶等于其循环分解中各循环长度的最小公倍数。［用练习 10 和第 1 节练习 24。］
  \item 证明若 $n\geq m$，则 $S_n$ 中 $m$-循环的个数为
  \[
    \frac{n(n-1)(n-2)\cdots(n-m+1)}{m}.
  \]
  ［数一数形成 $m$-循环的方法数，再除以一个特定 $m$-循环的表示数。］
  \item 证明若 $n\geq4$，则 $S_n$ 中可写成两个不交 $2$-循环之积的置换个数为
  \[
    \frac{n(n-1)(n-2)(n-3)}{8}.
  \]
  \item 求所有整数 $n$，使得 $S_5$ 含有阶为 $n$ 的元素。［用练习 15。］
  \item 求所有整数 $n$，使得 $S_7$ 含有阶为 $n$ 的元素。［用练习 15。］
  \item 找出 $S_3$ 的一组生成元与关系。
\end{enumerate}

\section{矩阵群}

本节引入系数来自域的矩阵群的概念。这样一族群将在第一部分中作为说明性例子使用，并会在关于向量空间的章节中更详细地研究。

域是能够进行所有算术运算 $+$、$-$、$\times$ 和 $/$（除以非零元素）的“最小”数学结构；特别地，每个非零元素都必须有乘法逆元。我们稍后会更深入地研究域；在本书这一部分中，我们遇到的域 $F$ 只会是 $\QQ$、$\RR$ 和 $\ZZ/p\ZZ$，其中 $p$ 是素数。例子 $\ZZ/p\ZZ$ 是一个有限域；为强调它是域，我们记作 $\FF_p$。为完整起见，这里给出域的精确定义。

\noindent\textbf{定义.}
\begin{enumerate}[label=(\arabic*)]
  \item \emph{域}是一个集合 $F$，连同 $F$ 上两个二元运算 $+$ 与 $\cdot$，使得 $(F,+)$ 是阿贝尔群（其单位元记为 $0$），$(F-\{0\},\cdot)$ 也是阿贝尔群，并且满足分配律
  \[
    a\cdot(b+c)=(a\cdot b)+(a\cdot c),\qquad\text{对所有 }a,b,c\in F.
  \]
  \item 对任意域 $F$，令 $F^\times=F-\{0\}$。
\end{enumerate}

当标量来自 $\RR$ 时，向量空间理论、矩阵与线性变换理论以及行列式理论中的所有结论，在作适当修改后，当标量来自任意域 $F$ 时也都成立。当我们在第一部分中使用这些理论时，会明确说明我们假设了哪些关于域的事实。

对每个 $n\in\ZZ^+$，令 $\GL_n(F)$ 为所有元素来自 $F$ 且行列式非零的 $n\times n$ 矩阵的集合，即
\[
  \GL_n(F)=\{A\mid A\text{ 是元素来自 }F\text{ 的 }n\times n\text{ 矩阵且 }\det(A)\neq0\}.
\]
任意元素来自 $F$ 的矩阵 $A$ 的行列式可用 $F=\RR$ 时的同样公式计算。对任意 $n\times n$ 矩阵 $A$ 与 $B$，令 $AB$ 为按 $F=\RR$ 时同样规则计算的矩阵积。这个乘法是结合的。此外，由 $\det(AB)=\det(A)\det(B)$，若 $\det(A)\neq0$ 且 $\det(B)\neq0$，则 $\det(AB)\neq0$，所以 $\GL_n(F)$ 对矩阵乘法封闭。进一步地，$\det(A)\neq0$ 当且仅当 $A$ 有矩阵逆（并且这个逆可由 $F=\RR$ 时使用的同样伴随矩阵公式计算），所以每个 $A\in\GL_n(F)$ 都有逆 $A^{-1}\in\GL_n(F)$：
\[
  AA^{-1}=A^{-1}A=I,
\]
其中 $I$ 是 $n\times n$ 单位矩阵。因此，$\GL_n(F)$ 在矩阵乘法下是群，称为\emph{$n$ 次一般线性群}。

以下结果将在第三部分中证明，这里为方便先记录：
\begin{enumerate}[label=(\arabic*)]
  \item 若 $F$ 是域且 $|F|<\infty$，则 $|F|=p^m$，其中 $p$ 为某个素数，$m$ 为某个整数；
  \item 若 $|F|=q<\infty$，则
  \[
    |\GL_n(F)|=(q^n-1)(q^n-q)(q^n-q^2)\cdots(q^n-q^{n-1}).
  \]
\end{enumerate}

\subsection*{练习}

设 $F$ 是域，$n\in\ZZ^+$。
\begin{enumerate}[label=\arabic*.]
  \item 证明 $|\GL_2(\FF_2)|=6$。
  \item 写出 $\GL_2(\FF_2)$ 的所有元素，并计算每个元素的阶。
  \item 证明 $\GL_2(\FF_2)$ 是非阿贝尔的。
  \item 证明若 $n$ 不是素数，则 $\ZZ/n\ZZ$ 不是域。
  \item 证明 $\GL_n(F)$ 是有限群，当且仅当 $F$ 含有有限多个元素。
  \item 若 $|F|=q$ 有限，证明 $|\GL_n(F)|<q^{n^2}$。
  \item 设 $p$ 是素数。证明 $\GL_2(\FF_p)$ 的阶为 $p^4-p^3-p^2+p$（不要直接引用本节中的阶公式）。［从 $\FF_p$ 上所有 $2\times2$ 矩阵总数中减去不可逆矩阵的个数。可以使用这样一个事实：$2\times2$ 矩阵不可逆，当且仅当一行是另一行的倍数。］
  \item 证明对任意 $n\geq2$ 和任意 $F$，$\GL_n(F)$ 都是非阿贝尔的。
  \item 证明实数项 $2\times2$ 矩阵的矩阵乘法是结合的。
  \item 令
  \[
    G=\left\{\begin{pmatrix}a&b\\0&c\end{pmatrix}\ \middle|\ a,b,c\in\RR,\ a\neq0,\ c\neq0\right\}.
  \]
  \begin{enumerate}[label=(\alph*)]
    \item 计算
    \[
      \begin{pmatrix}a_1&b_1\\0&c_1\end{pmatrix}
      \begin{pmatrix}a_2&b_2\\0&c_2\end{pmatrix}
    \]
    以说明 $G$ 对矩阵乘法封闭。
    \item 求矩阵 $\begin{pmatrix}a&b\\0&c\end{pmatrix}$ 的逆，并推出 $G$ 对取逆封闭。
    \item 推出 $G$ 是 $\GL_2(\RR)$ 的子群（参见第 1 节练习 26）。
    \item 证明 $G$ 中两个对角元相等（即 $a=c$）的元素组成的集合也是 $\GL_2(\RR)$ 的子群。
  \end{enumerate}
\end{enumerate}

下一个练习引入域 $F$ 上的海森堡群，并发展它的一些基本性质。当 $F=\RR$ 时，这个群在量子力学和信号理论中起重要作用，它给出了海森堡不确定性原理的一个群论解释（归功于 H. Weyl）。还要注意，海森堡群可以更一般地定义——例如，元素可取在 $\ZZ$ 中。
\begin{enumerate}[label=\arabic*.,resume]
  \item 令
  \[
    H(F)=\left\{\begin{pmatrix}1&a&b\\0&1&c\\0&0&1\end{pmatrix}\ \middle|\ a,b,c\in F\right\},
  \]
  称为 $F$ 上的\emph{海森堡群}。令
  \[
    X=\begin{pmatrix}1&a&b\\0&1&c\\0&0&1\end{pmatrix},\qquad
    Y=\begin{pmatrix}1&d&e\\0&1&f\\0&0&1\end{pmatrix}
  \]
  为 $H(F)$ 的元素。
  \begin{enumerate}[label=(\alph*)]
    \item 计算矩阵积 $XY$，并推出 $H(F)$ 对矩阵乘法封闭。给出显式矩阵，使得 $XY\neq YX$（从而 $H(F)$ 总是非阿贝尔的）。
    \item 求矩阵逆 $X^{-1}$ 的显式公式，并推出 $H(F)$ 对取逆封闭。
    \item 证明 $H(F)$ 的结合律，并推出 $H(F)$ 是阶为 $|F|^3$ 的群。（不要假设矩阵乘法结合。）
    \item 求有限群 $H(\ZZ/2\ZZ)$ 中每个元素的阶。
    \item 证明群 $H(\RR)$ 的每个非单位元素都是无限阶的。
  \end{enumerate}
\end{enumerate}

\section{四元数群}

\emph{四元数群} $Q_8$ 定义为
\[
  Q_8=\{1,-1,i,-i,j,-j,k,-k\},
\]
乘法按如下规则计算：
\begin{align*}
  1\cdot a&=a\cdot1=a, &&\text{对所有 }a\in Q_8,\\
  (-1)\cdot(-1)&=1,\qquad (-1)\cdot a=a\cdot(-1)=-a, &&\text{对所有 }a\in Q_8,\\
  i\cdot i&=j\cdot j=k\cdot k=-1,\\
  i\cdot j&=k, & j\cdot i&=-k,\\
  j\cdot k&=i, & k\cdot j&=-i,\\
  k\cdot i&=j, & i\cdot k&=-j.
\end{align*}
照常，今后把 $a\cdot b$ 写作 $ab$。检验结合律相当繁琐（我们稍后会用较少计算的方法证明它），但其他公理容易检验。注意，$Q_8$ 是阶为 $8$ 的非阿贝尔群。

\subsection*{练习}
\begin{enumerate}[label=\arabic*.]
  \item 计算 $Q_8$ 中每个元素的阶。
  \item 写出 $S_3$、$D_8$ 和 $Q_8$ 的群表。
  \item 找出 $Q_8$ 的一组生成元与关系。
\end{enumerate}


\section{同态与同构}

本节精确说明两个群何时“看起来相同”，也就是说，何时具有完全相同的群论结构。这就是两个群之间同构的概念。我们先定义同态的概念，稍后还会对此作大量讨论。

\noindent\textbf{定义.} 设 $(G,*)$ 和 $(H,\circ)$ 是群。若映射 $\varphi:G\to H$ 满足
\[
  \varphi(x*y)=\varphi(x)\circ\varphi(y),\qquad\text{对所有 }x,y\in G,
\]
则称 $\varphi$ 是一个\emph{同态}。

当 $G$ 和 $H$ 的群运算没有显式写出时，同态条件就简写为
\[
  \varphi(xy)=\varphi(x)\varphi(y),
\]
但必须牢记，左边的乘积 $xy$ 是在 $G$ 中计算的，而右边的乘积 $\varphi(x)\varphi(y)$ 是在 $H$ 中计算的。直观地说，若映射 $\varphi$ 保持其定义域与陪域的群结构，则它是同态。

\noindent\textbf{定义.} 映射 $\varphi:G\to H$ 称为一个\emph{同构}，并且称 $G$ 与 $H$ \emph{同构}或具有相同的\emph{同构型}，记作 $G\cong H$，若
\begin{enumerate}[label=(\arabic*)]
  \item $\varphi$ 是同态（即 $\varphi(xy)=\varphi(x)\varphi(y)$），并且
  \item $\varphi$ 是双射。
\end{enumerate}

换言之，群 $G$ 与 $H$ 同构，是指存在一个保持群运算的双射从 $G$ 到 $H$。直观地说，除元素名称和运算符号可能写得不同外，$G$ 与 $H$ 是同一个群。因此，$G$ 所具有的任何只依赖于 $G$ 的群结构的性质（也就是可由群公理推出的性质——例如群的交换性）也在 $H$ 中成立。注意，这从形式上说明了为什么可以把所有群运算都写作 $\cdot$：改变运算符号并不改变同构型。

\paragraph{例}
\begin{enumerate}[label=(\arabic*)]
  \item 对任意群 $G$，$G\cong G$。恒等映射给出了一个明显的同构，但一般并不是从 $G$ 到自身的唯一同构。更一般地，令 $\mathcal{G}$ 为任意非空群族。容易检验关系 $\cong$ 是 $\mathcal{G}$ 上的等价关系，其等价类称为\emph{同构类}。这也解释了“同构”定义中略显对称的措辞。
  \item 指数映射 $\exp:\RR\to\RR^+$，定义为 $\exp(x)=e^x$（其中 $e$ 是自然对数的底），是从 $(\RR,+)$ 到 $(\RR^+,\times)$ 的同构。$\exp$ 是双射，因为它有逆函数（即 $\log_e$）；且 $\exp$ 保持群运算，因为 $e^{x+y}=e^xe^y$。在这个例子中，元素和运算都不同，但两个群同构；也就是说，作为群它们有相同的结构。
  \item 本例说明，对称群的同构型只依赖于被置换的底层集合的基数。令 $\Delta$ 与 $\Omega$ 是非空集合。若 $|\Delta|=|\Omega|$，则对称群 $S_\Delta$ 与 $S_\Omega$ 同构。直观上，既然 $|\Delta|=|\Omega|$，存在一个从 $\Delta$ 到 $\Omega$ 的双射 $\theta$。把 $\Delta$ 与 $\Omega$ 的元素通过 $\theta$ 粘在一起，即每个 $x\in\Delta$ 与 $\theta(x)\in\Omega$ 粘在一起。为了得到映射 $\Phi:S_\Delta\to S_\Omega$，令 $\sigma\in S_\Delta$ 为 $\Delta$ 的一个置换，并令 $\Phi(\sigma)$ 为 $\Omega$ 的一个置换，它以 $\sigma$ 移动相应被粘合的 $\Delta$ 元素的方式来移动 $\Omega$ 的元素；也就是说，若对某些 $x,y\in\Delta$ 有 $\sigma(x)=y$，则在 $\Omega$ 中令 $\Phi(\sigma)(\theta(x))=\theta(y)$。由于集合双射 $\theta$ 有逆，容易检验对称群之间的这个映射也有逆。该映射的精确定义及其作为同构所需性质的直接（虽有些繁琐）验证留给后面的练习。

  反过来，若 $S_\Delta\cong S_\Omega$，则 $|\Delta|=|\Omega|$；我们只在底层集合有限时证明这一点（当 $\Delta$ 与 $\Omega$ 都是无限集时，证明更难，将在第 4 章作为练习给出）。任意两个群 $G$ 与 $H$ 的同构，首先就是二者之间的双射，因此同构的必要条件是 $|S_\Delta|=|S_\Omega|$。当 $\Delta$ 是阶为 $n$ 的有限集时，$|S_\Delta|=n!$；虽然我们只为 $S_n$ 证明过这一点，但同样的论证适用于 $S_\Delta$。类似地，若 $\Omega$ 是阶为 $m$ 的有限集，则 $|S_\Omega|=m!$。因此若 $S_\Delta$ 与 $S_\Omega$ 同构，则 $n!=m!$，所以 $m=n$，即 $|\Delta|=|\Omega|$。
\end{enumerate}

全文还会出现更多同构的例子。当我们研究不同结构（环、域、向量空间等）时，也会为相应结构之间制定对应的同构概念。数学中的一个核心问题是确定一个结构的哪些性质能够指定它的同构型（即证明：若 $G$ 是具有某种结构的对象（例如群），且 $G$ 具有性质 $P$，则任何其他同类结构对象（群）$X$ 只要具有性质 $P$，就同构于 $G$）。这类定理称为\emph{分类定理}。例如，我们将证明
\[
  \text{任意阶为 }6\text{ 的非阿贝尔群都同构于 }S_3.
\]
（这里 $G$ 是群 $S_3$，$P$ 是性质“非阿贝尔且阶为 $6$”。）由这个分类定理，我们无需寻找这些群之间的显式映射，就得到 $D_6\cong S_3$ 以及 $\GL_2(\FF_2)\cong S_3$。注意，并非任意阶为 $6$ 的群都同构于 $S_3$。事实上，我们将证明在同构意义下，阶为 $6$ 的群恰有两个：$S_3$ 与 $\ZZ/6\ZZ$（也就是说，任意阶为 $6$ 的群都同构于这两个群之一，且 $S_3$ 不同构于 $\ZZ/6\ZZ$）。注意，结论不那么具体（有两种可能类型），但假设更容易检验（只需检验阶是否为 $6$）。后一类结果也称为分类。一般来说，判断两个群（或其他数学对象）是否同构，即使在具体例子中也是微妙而困难的——构造一个保持群运算的显式映射，或证明不存在这样的映射，除极小的情形外通常在计算上不可行；这一点在没有进一步理论而试图证明上述阶 $6$ 群分类时已经有所体现。

有时容易看出两个给定群不同构。例如，下面练习断言若 $\varphi:G\to H$ 是同构，则特别地：
\begin{enumerate}[label=(\alph*)]
  \item $|G|=|H|$；
  \item $G$ 是阿贝尔群，当且仅当 $H$ 是阿贝尔群；
  \item 对所有 $x\in G$，$|x|=|\varphi(x)|$。
\end{enumerate}
因此，$S_3$ 与 $\ZZ/6\ZZ$ 不同构（如上所述），因为一个是阿贝尔的，另一个不是。还有，$(\RR-\{0\},\times)$ 与 $(\RR,+)$ 不可能同构，因为在 $(\RR-\{0\},\times)$ 中元素 $-1$ 的阶为 $2$，而 $(\RR,+)$ 中没有阶为 $2$ 的元素，这与 (c) 矛盾。

最后，记录一个稍后（讨论自由群时）将证明的非常有用的事实，它处理由生成元与关系给出的两个群之间同态与同构的问题：

设 $G$ 是一个有限群，阶为 $n$，我们有它的一个表示，并令 $S=\{s_1,\ldots,s_m\}$ 为生成元。设 $H$ 是另一个群，且 $\{r_1,\ldots,r_m\}$ 是 $H$ 中元素。假设 $G$ 中由 $s_i$ 满足的任一关系，在把每个 $s_i$ 换成 $r_i$ 后也在 $H$ 中成立。则存在唯一同态 $\varphi:G\to H$，把 $s_i$ 映到 $r_i$。若我们有 $G$ 的一个表示，则只需检验该表示所列出的关系（因为按表示的定义，所有关系都能由表示中给出的关系推出）。若 $H$ 由元素 $\{r_1,\ldots,r_m\}$ 生成，则 $\varphi$ 是满射（任意 $r_i$ 的乘积都是相应 $s_i$ 乘积的像）。若此外 $H$ 与 $G$ 有相同的有限阶，则任意满射必为单射，因此 $\varphi$ 是同构：$G\cong H$。直观地说，我们可以把 $G$ 的生成元映到 $H$ 的任意元素上，并在 $G$ 中的关系仍被满足的条件下得到一个同态。

读者可能已经熟悉向量空间中的相应说法。设 $V$ 是一个 $n$ 维有限维向量空间，基为 $S$，$W$ 是另一个向量空间。则可以通过把 $S$ 中元素映到 $W$ 中任意向量来指定从 $V$ 到 $W$ 的一个线性变换（这里没有关系需要满足）。若 $W$ 也维数为 $n$，并且在 $W$ 中选定的这些向量张成 $W$（因而是 $W$ 的一组基），则该线性变换可逆（是向量空间同构）。

\paragraph{例}
\begin{enumerate}[label=(\arabic*)]
  \item 回忆 $D_{2n}=\langle r,s\mid r^n=s^2=1,\ sr=r^{-1}s\rangle$。假设 $H$ 是一个含有元素 $a,b$ 的群，并满足 $a^n=1$、$b^2=1$ 和 $ba=a^{-1}b$。则存在一个从 $D_{2n}$ 到 $H$ 的同态，把 $r$ 映到 $a$，把 $s$ 映到 $b$。例如，令 $k$ 是整除 $n$ 的整数且 $k\geq3$，并令
  \[
    D_{2k}=\langle r_1,s_1\mid r_1^k=s_1^2=1,\ s_1r_1=r_1^{-1}s_1\rangle.
  \]
  定义
  \[
    \varphi:D_{2n}\to D_{2k}\quad\text{由}\quad \varphi(r)=r_1,
    \quad \varphi(s)=s_1.
  \]
  若写 $n=km$，则因为 $r_1^k=1$，也有 $r_1^n=(r_1^k)^m=1$。于是 $D_{2n}$ 中由 $r,s$ 满足的三个关系也由 $D_{2k}$ 中的 $r_1,s_1$ 满足。因此 $\varphi$ 唯一地扩张为从 $D_{2n}$ 到 $D_{2k}$ 的同态。由于 $\{r_1,s_1\}$ 生成 $D_{2k}$，$\varphi$ 是满射。若 $k<n$，这个同态不是同构。
  \item 继续前例，令 $G=D_6$ 如上表示。检验在 $H=S_3$ 中，元素 $a=(1\ 2\ 3)$ 与 $b=(1\ 2)$ 满足关系：$a^3=1$、$b^2=1$ 且 $ba=a^{-1}b$。因此存在从 $D_6$ 到 $S_3$ 的同态，把 $r\mapsto a$，$s\mapsto b$。进一步可检验 $S_3$ 由 $a$ 与 $b$ 生成，所以此同态是满射。由于 $D_6$ 与 $S_3$ 的阶都为 $6$，这个同态是同构：$D_6\cong S_3$。
\end{enumerate}

注意，上述例子中的元素 $a$ 不必有阶 $n$（即 $n$ 不必是使 $a$ 的幂等于 $H$ 中单位元的最小正指数），同样，$b$ 也不必有阶 $2$（例如若 $a=a^{-1}$，则 $b$ 完全可以是单位元）。这使我们能更容易地构造同态，也符合这样的思想：群 $G$ 的生成元与关系构成了 $G$ 的群结构的一套完整数据。

\subsection*{练习}
设 $G$ 与 $H$ 是群。
\begin{enumerate}[label=\arabic*.]
  \item 设 $\varphi:G\to H$ 是同态。
  \begin{enumerate}[label=(\alph*)]
    \item 证明对所有 $n\in\ZZ^+$，$\varphi(x^n)=\varphi(x)^n$。
    \item 对 $n=-1$ 证明 (a)，并推出对所有 $n\in\ZZ$，$\varphi(x^n)=\varphi(x)^n$。
  \end{enumerate}
  \item 若 $\varphi:G\to H$ 是同构，证明对所有 $x\in G$，$|\varphi(x)|=|x|$。推出任意两个同构群对每个 $n\in\ZZ^+$ 都有相同个数的阶为 $n$ 的元素。若 $\varphi$ 仅假设为同态，结论是否仍成立？
  \item 若 $\varphi:G\to H$ 是同构，证明 $G$ 是阿贝尔群当且仅当 $H$ 是阿贝尔群。若 $\varphi:G\to H$ 是同态，还需要对 $\varphi$ 加上什么条件（若需要），才能保证若 $G$ 阿贝尔，则 $H$ 也阿贝尔？
  \item 证明乘法群 $\RR-\{0\}$ 与 $\CC-\{0\}$ 不同构。
  \item 证明加法群 $\RR$ 与 $\QQ$ 不同构。
  \item 证明加法群 $\ZZ$ 与 $\QQ$ 不同构。
  \item 证明 $D_8$ 与 $Q_8$ 不同构。
  \item 证明若 $n\neq m$，则 $S_n$ 与 $S_m$ 不同构。
  \item 证明 $D_{24}$ 与 $S_4$ 不同构。
  \item 按如下方式补全证明：若 $|\Delta|=|\Omega|$，则对称群 $S_\Delta$ 与 $S_\Omega$ 同构。令 $\theta:\Delta\to\Omega$ 为双射，定义
  \[
    \varphi:S_\Delta\to S_\Omega,
    \qquad \varphi(\sigma)=\theta\circ\sigma\circ\theta^{-1},
    \quad \sigma\in S_\Delta.
  \]
  证明：
  \begin{enumerate}[label=(\alph*)]
    \item $\varphi$ 是良定的，即若 $\sigma$ 是 $\Delta$ 的置换，则 $\theta\circ\sigma\circ\theta^{-1}$ 是 $\Omega$ 的置换；
    \item $\varphi$ 是从 $S_\Delta$ 到 $S_\Omega$ 上的双射；［找出 $\varphi$ 的双侧逆。］
    \item $\varphi$ 是同态，即 $\varphi(\sigma\circ\tau)=\varphi(\sigma)\circ\varphi(\tau)$。
  \end{enumerate}
  注意这与矩阵的换基或相似变换的相似性（稍后我们会看到它们之间的联系）。
  \item 设 $A$ 与 $B$ 是群。证明 $A\times B\cong B\times A$。
  \item 设 $A,B,C$ 是群，令 $G=A\times B$，$H=B\times C$。证明 $G\times C\cong A\times H$。
  \item 设 $G,H$ 是群，$\varphi:G\to H$ 是同态。证明 $\varphi$ 的像 $\varphi(G)$ 是 $H$ 的子群（参见第 1 节练习 26）。证明若 $\varphi$ 是单射，则 $G\cong\varphi(G)$。
  \item 设 $G,H$ 是群，$\varphi:G\to H$ 是同态。定义 $\varphi$ 的\emph{核}为 $\{g\in G\mid \varphi(g)=1_H\}$（即 $G$ 中映到 $H$ 的单位元的元素集合，也就是 $H$ 的单位元上方的纤维）。证明 $\varphi$ 的核是 $G$ 的子群（参见第 1 节练习 26）。证明 $\varphi$ 是单射当且仅当 $\varphi$ 的核是 $G$ 的单位子群。
  \item 定义映射 $\pi:\RR^2\to\RR$ 为 $\pi((x,y))=x$。证明 $\pi$ 是同态，并求 $\pi$ 的核（参见练习 14）。
  \item 设 $A$ 与 $B$ 是群，$G=A\times B$ 为它们的直积。证明由 $\pi_1((a,b))=a$ 与 $\pi_2((a,b))=b$ 定义的映射 $\pi_1:G\to A$ 和 $\pi_2:G\to B$ 是同态，并求它们的核（参见练习 14）。
  \item 设 $G$ 是任意群。证明从 $G$ 到自身的映射 $g\mapsto g^{-1}$ 是同态，当且仅当 $G$ 是阿贝尔群。
  \item 设 $G$ 是任意群。证明从 $G$ 到自身的映射 $g\mapsto g^2$ 是同态，当且仅当 $G$ 是阿贝尔群。
  \item 令 $G=\{z\in\CC\mid z^n=1\text{ 对某个 }n\in\ZZ^+\}$。证明对任意固定整数 $k>1$，由 $z\mapsto z^k$ 定义的从 $G$ 到自身的映射是满同态，但不是同构。
  \item 设 $G$ 是群，$\Aut(G)$ 为从 $G$ 到自身的所有同构组成的集合。证明 $\Aut(G)$ 在函数复合下是群（称为 $G$ 的自同构群，$\Aut(G)$ 的元素称为 $G$ 的自同构）。
  \item 证明对每个固定的非零 $k\in\QQ$，由 $q\mapsto kq$ 定义的从 $\QQ$ 到自身的映射是 $\QQ$ 的自同构（参见练习 20）。
  \item 设 $A$ 是阿贝尔群，并固定某个 $k\in\ZZ$。证明映射 $a\mapsto a^k$ 是从 $A$ 到自身的同态。若 $k=-1$，证明这个同态是同构（即 $A$ 的自同构）。
  \item 设 $G$ 是有限群，并且有一个自同构 $\alpha$（参见练习 20）满足 $\alpha(g)=g$ 当且仅当 $g=1$。若 $\alpha^2$ 是从 $G$ 到 $G$ 的恒等映射，证明 $G$ 是阿贝尔群（这样的自同构 $\alpha$ 称为阶为 $2$ 的无不动点自同构）。［证明 $G$ 的每个元素都可写成 $x^{-1}\alpha(x)$ 的形式，并对这样的表达式作用 $\alpha$。］
  \item 设 $G$ 是有限群，且 $x,y$ 是 $G$ 中两个不同的阶为 $2$ 的元素，它们生成 $G$。证明 $G\cong D_{2n}$，其中 $n=|xy|$。［见第 2 节练习 6。］
  \item 设 $n\in\ZZ^+$，$r,s$ 是 $D_{2n}$ 的通常生成元，令 $\theta=2\pi/n$。
  \begin{enumerate}[label=(\alph*)]
    \item 证明矩阵
    \[
      \begin{pmatrix}\cos\theta&-\sin\theta\\ \sin\theta&\cos\theta\end{pmatrix}
    \]
    是一个线性变换的矩阵，该线性变换把 $x,y$ 平面绕原点逆时针旋转 $\theta$ 弧度。
    \item 证明在生成元上由
    \[
      \varphi(r)=\begin{pmatrix}\cos\theta&-\sin\theta\\ \sin\theta&\cos\theta\end{pmatrix},
      \qquad
      \varphi(s)=\begin{pmatrix}0&1\\1&0\end{pmatrix}
    \]
    定义的映射 $\varphi:D_{2n}\to\GL_2(\RR)$ 可扩张为从 $D_{2n}$ 到 $\GL_2(\RR)$ 的同态。
    \item 证明 (b) 中的同态 $\varphi$ 是单射。
  \end{enumerate}
  \item 设 $i$ 和 $j$ 是第 5 节中描述的 $Q_8$ 的生成元。证明从 $Q_8$ 到 $\GL_2(\CC)$ 的映射 $\varphi$ 在生成元上由
  \[
    \varphi(i)=\begin{pmatrix}\sqrt{-1}&0\\0&-\sqrt{-1}\end{pmatrix},
    \qquad
    \varphi(j)=\begin{pmatrix}0&-1\\1&0\end{pmatrix}
  \]
  定义时，可扩张为一个同态。证明 $\varphi$ 是单射。
\end{enumerate}


\section{群作用}

本节给出群作用在集合上的精确定义，并给出一些例子。群作用将是一个强有力的工具，我们将用它来证明抽象群的定理，并揭示具体例子的结构。此外，“作用”这一概念会作为研究代数对象的一种方法反复出现：通过考察该对象如何作用在其他结构上来研究它。

\noindent\textbf{定义.} 群 $G$ 在集合 $A$ 上的一个\emph{群作用}，是一个从 $G\times A$ 到 $A$ 的映射（对所有 $g\in G$ 与 $a\in A$，写作 $g\cdot a$），满足下列性质：
\begin{enumerate}[label=(\arabic*)]
  \item 对所有 $g_1,g_2\in G$、$a\in A$，有 $g_1\cdot(g_2\cdot a)=(g_1g_2)\cdot a$；
  \item 对所有 $a\in A$，有 $1\cdot a=a$。
\end{enumerate}

我们很快会不那么形式化，并说 $G$ 是一个作用在集合 $A$ 上的群。当不致与例如群运算混淆时，表达式 $g\cdot a$ 通常简写为 $ga$（请记住，这里的 $\cdot$ 不是 $A$ 上的二元运算，并且 $ga$ 总是 $A$ 的一个元素）。注意，在性质 (1) 的左边，$g_2\cdot a$ 是 $A$ 的元素，所以用 $g_1$ 作用在它上面是有意义的；在该等式右边，乘积 $g_1g_2$ 是在 $G$ 中计算的，所得群元素再作用在集合元素 $a$ 上。

在给出群作用的例子之前，先作一些观察。设群 $G$ 作用在集合 $A$ 上。对每个固定的 $g\in G$，得到映射 $\sigma_g:A\to A$，定义为
\[
  \sigma_g(a)=g\cdot a.
\]
我们证明两个重要事实：
\begin{enumerate}[label=(\roman*)]
  \item 对每个固定的 $g\in G$，$\sigma_g$ 是 $A$ 的一个置换；
  \item 由 $g\mapsto\sigma_g$ 定义的从 $G$ 到 $S_A$ 的映射是同态。
\end{enumerate}

为说明 $\sigma_g$ 是 $A$ 的置换，我们证明它作为 $A$ 到 $A$ 的集合映射有一个双侧逆，即 $\sigma_{g^{-1}}$（于是由第 0.1 节命题 1，它是置换）。对所有 $a\in A$，
\begin{align*}
  (\sigma_{g^{-1}}\circ\sigma_g)(a)
  &=\sigma_{g^{-1}}(\sigma_g(a)) &&\text{（函数复合的定义）}\\
  &=g^{-1}\cdot(g\cdot a) &&\text{（$\sigma_{g^{-1}}$ 与 $\sigma_g$ 的定义）}\\
  &=(g^{-1}g)\cdot a &&\text{（作用的性质 (1)）}\\
  &=1\cdot a=a &&\text{（作用的性质 (2)）。}
\end{align*}
这证明 $\sigma_{g^{-1}}\circ\sigma_g$ 是 $A$ 到 $A$ 的恒等映射。由于 $g$ 是任意的，交换 $g$ 与 $g^{-1}$ 的角色可知 $\sigma_g\circ\sigma_{g^{-1}}$ 也是 $A$ 上的恒等映射。因此 $\sigma_g$ 有双侧逆，故是 $A$ 的置换。

为检验上面的断言 (ii)，定义 $\varphi:G\to S_A$ 为 $\varphi(g)=\sigma_g$。注意 (i) 说明 $\sigma_g$ 的确是 $S_A$ 的元素。为了证明 $\varphi$ 是同态，必须证明
\[
  \varphi(g_1g_2)=\varphi(g_1)\circ\varphi(g_2)
\]
（回忆 $S_A$ 在函数复合下成群）。置换 $\varphi(g_1g_2)$ 与 $\varphi(g_1)\circ\varphi(g_2)$ 相等，当且仅当它们在每个 $a\in A$ 上的取值相同。对所有 $a\in A$，
\begin{align*}
  \varphi(g_1g_2)(a)
  &=\sigma_{g_1g_2}(a) &&\text{（$\varphi$ 的定义）}\\
  &=(g_1g_2)\cdot a &&\text{（$\sigma_{g_1g_2}$ 的定义）}\\
  &=g_1\cdot(g_2\cdot a) &&\text{（作用的性质 (1)）}\\
  &=\sigma_{g_1}(\sigma_{g_2}(a)) &&\text{（$\sigma_{g_1}$ 与 $\sigma_{g_2}$ 的定义）}\\
  &=(\varphi(g_1)\circ\varphi(g_2))(a) &&\text{（$\varphi$ 的定义）。}
\end{align*}
这证明了断言 (ii)。

直观地说，群 $G$ 在集合 $A$ 上的作用，就是 $G$ 中每个元素 $g$ 都以一种与 $G$ 的群运算相容的方式作为 $A$ 上的置换起作用；上面的断言 (i) 和 (ii) 使这一点精确化。上述从 $G$ 到 $S_A$ 的同态称为给定作用所伴随的\emph{置换表示}。容易看出，这个过程可以反过来：若 $\varphi:G\to S_A$ 是从群 $G$ 到集合 $A$ 上的对称群的任意同态，则由
\[
  g\cdot a=\varphi(g)(a),\qquad g\in G,
  \quad a\in A,
\]
定义的从 $G\times A$ 到 $A$ 的映射满足 $G$ 在 $A$ 上的群作用性质。因此，群 $G$ 在集合 $A$ 上的作用与从 $G$ 到对称群 $S_A$ 的同态一一对应（也就是说，本质上是同一概念，只是用不同术语表述）。还应注意，作用的这个定义更准确地说应称为\emph{左作用}，因为群元素写在集合元素的左边。类似地也可以定义右作用。

\paragraph{例}
设 $G$ 是群，$A$ 是非空集合。下面每个例子中，对作用性质 (1) 与 (2) 的检验都留作练习。
\begin{enumerate}[label=(\arabic*)]
  \item 令 $ga=a$，对所有 $g\in G$ 与 $a\in A$。群作用的性质 (1) 和 (2) 立即成立。这个作用称为\emph{平凡作用}，并称 $G$ 在 $A$ 上\emph{平凡地作用}。注意，$G$ 的不同元素诱导 $A$ 上同一个置换（在此情形中为恒等置换）。相应的置换表示 $G\to S_A$ 是把每个 $G$ 中元素都映到恒等置换的平凡同态。

  若 $G$ 作用在集合 $B$ 上，且 $G$ 的不同元素诱导 $B$ 上不同的置换，则称该作用是\emph{忠实的}。因此，忠实作用正是相应置换表示为单射的作用。

  $G$ 在 $B$ 上作用的\emph{核}定义为
  \[
    \{g\in G\mid gb=b\text{ 对所有 }b\in B\},
  \]
  即固定 $B$ 中所有元素的 $G$ 中元素。对平凡作用，作用的核为整个 $G$；当 $|G|>1$ 时，该作用不是忠实的。

  \item 域 $F$ 上向量空间 $V$ 的公理包括：乘法群 $F^\times$ 作用在集合 $V$ 上。因此，向量空间是域的乘法群作用的熟悉例子，而且这里还有更多结构（特别是 $V$ 必须是阿贝尔群）可加以利用。在特殊情形 $V=\RR^n$ 且 $F=\RR$ 时，该作用由
  \[
    a(r_1,r_2,\ldots,r_n)=(ar_1,ar_2,\ldots,ar_n)
  \]
  给出，其中 $a\in\RR$，$(r_1,r_2,\ldots,r_n)\in\RR^n$，$ar_i$ 只是两个实数的乘法。

  \item 对任意非空集合 $A$，对称群 $S_A$ 通过 $\sigma\cdot a=\sigma(a)$（对所有 $\sigma\in S_A$，$a\in A$）作用在 $A$ 上。相应的置换表示就是从 $S_A$ 到自身的恒等映射。

  \item 若固定正 $n$ 边形顶点的一个标号，则 $D_{2n}$ 的每个元素 $\alpha$ 都根据该对称如何置换对应顶点，给出 $\{1,2,\ldots,n\}$ 的一个置换 $\sigma_\alpha$。由
  \[
    (\alpha,i)\mapsto \sigma_\alpha(i)
  \]
  定义的从 $D_{2n}\times\{1,2,\ldots,n\}$ 到 $\{1,2,\ldots,n\}$ 的映射，给出 $D_{2n}$ 在 $\{1,2,\ldots,n\}$ 上的群作用。按照群作用的记号，我们现在可以不用形式化且笨重的 $\sigma_\alpha(i)$，而写作 $\alpha i$。注意这个作用是忠实的：正 $n$ 边形的不同对称诱导顶点上的不同置换。

  当 $n=3$ 时，$D_6$ 在三角形三个（带标号）顶点上的作用给出了从 $D_6$ 到 $S_3$ 的单同态。由于这两个群阶相同，该映射也必为满射，即为同构：$D_6\cong S_3$。这是前一节通过生成元与关系建立的同一事实的另一个证明。几何上它表示：三角形顶点的任意置换都是一个对称。对任意 $n\geq4$ 的 $n$ 边形，类似说法不成立（仅从阶数考虑，就不可能有 $D_{2n}$ 与 $S_n$ 同构）。

  \item 设 $G$ 是任意群，并令 $A=G$。定义从 $G\times A$ 到 $A$ 的映射为 $g\cdot a=ga$，对每个 $g\in G$ 与 $a\in A$，其中右侧的 $ga$ 是群 $G$ 中 $g$ 与 $a$ 的乘积。这样得到 $G$ 在自身上的一个群作用，每个固定的 $g\in G$ 都通过左乘置换 $G$ 中元素：
  \[
    g:a\mapsto ga\qquad\text{对所有 }a\in G.
  \]
  （若 $G$ 用加法记号，则得到 $a\mapsto g+a$，并称为左平移。）这个作用称为 $G$ 在自身上的\emph{左正则作用}。由消去律可知，这个作用是忠实的（请检验）。
\end{enumerate}

其他作用的例子见练习。

\subsection*{练习}
\begin{enumerate}[label=\arabic*.]
  \item 设 $F$ 是域。证明 $F$ 的非零元素乘法群（记为 $F^\times$）通过 $g\cdot a=ga$ 作用在集合 $F$ 上，其中 $g\in F^\times$、$a\in F$，而 $ga$ 是 $F$ 中两个域元素的通常乘积（请明确指出使用了域定义中的哪些公理）。
  \item 证明加法群 $\ZZ$ 通过 $z\cdot a=z+a$（对所有 $z,a\in\ZZ$）作用在自身上。
  \item 证明加法群 $\RR$ 通过 $r\cdot(x,y)=(x+ry,y)$ 作用在 $x,y$ 平面 $\RR\times\RR$ 上。
  \item 设群 $G$ 作用在集合 $A$ 上，并固定某个 $a\in A$。证明下列集合是 $G$ 的子群（参见第 1 节练习 26）：
  \begin{enumerate}[label=(\alph*)]
    \item 该作用的核；
    \item $\{g\in G\mid ga=a\}$——这个子群称为 $a$ 在 $G$ 中的\emph{稳定子}。
  \end{enumerate}
  \item 证明群 $G$ 在集合 $A$ 上作用的核，与相应置换表示 $G\to S_A$ 的核相同（参见第 6 节练习 14）。
  \item 证明群 $G$ 在集合 $A$ 上忠实作用，当且仅当该作用的核只由单位元组成。
  \item 证明本节例 2 中的作用是忠实的。
  \item 设 $A$ 是非空集合，$k$ 是正整数且 $k\leq |A|$。对称群 $S_A$ 通过
  \[
    \sigma\cdot\{a_1,\ldots,a_k\}=\{\sigma(a_1),\ldots,\sigma(a_k)\}
  \]
  作用在由 $A$ 的所有 $k$ 元子集组成的集合 $B$ 上。
  \begin{enumerate}[label=(\alph*)]
    \item 证明这是一个群作用。
    \item 明确描述元素 $(1\ 2)$ 和 $(1\ 2\ 3)$ 如何作用在 $\{1,2,3,4\}$ 的六个 $2$ 元子集上。
  \end{enumerate}
  \item 把前一题中的“$k$ 元子集”换成“有序 $k$ 元组”，完成同样的两个部分；其中对 $k$ 元组的作用按上式定义，但把集合花括号换成圆括号（注意，例如 $2$ 元组 $(1,2)$ 与 $(2,1)$ 不同，而集合 $\{1,2\}$ 与 $\{2,1\}$ 相同，所以被作用的集合不同）。
  \item 参考前两题，判断：
  \begin{enumerate}[label=(\alph*)]
    \item 对哪些 $k$，$S_n$ 在 $k$ 元子集上的作用是忠实的；
    \item 对哪些 $k$，$S_n$ 在有序 $k$ 元组上的作用是忠实的。
  \end{enumerate}
  \item 写出与 $D_8$ 作用在正方形顶点上相对应的 $S_4$ 中八个置换的循环分解（正方形顶点按第 2 节中的方式标号）。
  \item 假设 $n$ 为偶正整数，证明 $D_{2n}$ 作用在由正 $n$ 边形的相对顶点对组成的集合上。求这个作用的核（照常标号顶点）。
  \item 求左正则作用的核。
  \item 设 $G$ 是群并令 $A=G$。证明若 $G$ 非阿贝尔，则由 $g\cdot a=ag$（对所有 $g,a\in G$）定义的映射不满足 $G$ 在自身上的（左）群作用公理。
  \item 设 $G$ 是任意群并令 $A=G$。证明由 $g\cdot a=ag^{-1}$（对所有 $g,a\in G$）定义的映射满足 $G$ 在自身上的（左）群作用公理。
  \item 设 $G$ 是任意群并令 $A=G$。证明由 $g\cdot a=gag^{-1}$（对所有 $g,a\in G$）定义的映射满足（左）群作用公理（这个 $G$ 在自身上的作用称为\emph{共轭作用}）。
  \item 设 $G$ 是群，并令 $G$ 通过左共轭作用在自身上，所以每个 $g\in G$ 把 $G$ 映到 $G$：
  \[
    x\mapsto gxg^{-1}.
  \]
  对固定的 $g\in G$，证明由 $g$ 的共轭给出的映射是从 $G$ 到自身的同构（即 $G$ 的自同构——参见第 6 节练习 20）。推出对所有 $x\in G$，$x$ 与 $gxg^{-1}$ 有相同的阶；并且对 $G$ 的任意子集 $A$，有 $|A|=|gAg^{-1}|$（其中 $gAg^{-1}=\{gag^{-1}\mid a\in A\}$）。
  \item 设 $H$ 是作用在集合 $A$ 上的群。证明 $A$ 上由
  \[
    a\sim b\quad\text{当且仅当}\quad a=hb\text{ 对某个 }h\in H
  \]
  定义的关系 $\sim$ 是等价关系。（对每个 $x\in A$，$x$ 在 $\sim$ 下的等价类称为 $x$ 在 $H$ 的作用下的\emph{轨道}。$H$ 作用下的轨道划分集合 $A$。）
  \item 设 $H$ 是有限群 $G$ 的子群（参见第 1 节练习 26），并令 $H$ 通过左乘作用在 $G$ 上（这里 $A=G$）。设 $x\in G$，令 $O$ 为 $x$ 在 $H$ 的作用下的轨道。证明映射
  \[
    H\to O,\qquad h\mapsto hx
  \]
  是双射（因此所有轨道的基数都是 $|H|$）。由此和前一题推出拉格朗日定理：
  \[
    \text{若 }G\text{ 是有限群且 }H\text{ 是 }G\text{ 的子群，则 }|H|\text{ 整除 }|G|.
  \]
  \item 说明正四面体的刚体运动群同构于 $S_4$ 的一个子群（参见第 1 节练习 26）。
  \item 说明立方体的刚体运动群同构于 $S_4$。［这个群作用在四对相对顶点组成的集合上。］
  \item 说明正八面体的刚体运动群同构于 $S_4$ 的一个子群（参见第 1 节练习 26）。［这个群作用在四对相对面组成的集合上。］推出立方体与正八面体的刚体运动群同构。（这些群同构是因为这些立体互为“对偶”——见 H. Coxeter, \textit{Introduction to Geometry}, Wiley, 1961。稍后我们会看到正十二面体与正二十面体的刚体运动群也同构——这些立体同样互为对偶。）
  \item 解释为什么立方体的刚体运动群在三对相对面组成的集合上的作用不是忠实的。求该作用的核。
\end{enumerate}
